
\section{Middle primes I: blocks and spectral bounds for the twelve-label law}
\label{sec:middle-reference}

\paragraph{What this section proves and how to read it.}
The middle range is $p\equiv3\pmod4$ in
$[Q_3,P_3)=[19\,951\,586,167\cdot2^{40})$. This section partitions
it into prime blocks $\Blk$ with $\sum_{p\in\Blk}\log p\le570$
(grouped into $33$ consecutive intervals, called \emph{pieces} in
displays), proves degree-sum bounds converting per-prime constants
into block prices (sums $\sum q_{\Blk}^{2a}b_{\Blk}^{-d/2}$; ``block
charge'' below means such a sum), and proves spectral bounds
($-gp^{-1/2}I\preceq A_p\preceq(1+c^+/p)I$, called \emph{caps}) for
the twelve-label graph laws. The summand of weight $1$ in the five-type
completion of \cref{lem:front12:reference} is called the \emph{principal law} (\emph{root} in this part; the word has
nothing to do with square roots). Cut codes name mirror symmetries:
L20 is one chamber symmetry, expanded at first use. Inputs: Part~IV.
Outputs: blocks and caps used to build $\mu,\lambda$ in
\cref{sec:middle-frame}. On seven intervals from $2^{30}$ the
principal law is reweighted; \cref{sec:positive-roots} bounds it.
A first reading should follow the block definition and cap
statements; threshold arithmetic is reference material.

\emph{Conventions.} In this section $\vartheta_a$ is the chamber cut
$\vartheta_{0,a}$ of the first coordinate: the reflection
$u_0\mapsto2a-1-u_0$ with plus half $H_a=\{a,a+1,a+2\}$ (indices mod $6$),
$a\in\Z/6\Z$. The reserved cut $\iota$ flips the layer $\ell$ and has plus
half $V_+=\{\ell=0\}$. For a prime $p$ we write $r=p^{-1/2}$ and $q=1/p$; on a
block $\Blk$ with least prime $b=b_{\Blk}$, $r=b^{-1/2}$ and $q=1/b$, while
$q_{\Blk}=\prod_{p\in\Blk}p$. The order $\succeq_c$ of cut matrices is that of
\cref{def:cc:order}. Every middle block has $\log q_{\Blk}\le L=570$; the
local data of a piece $[X,Y)$ are evaluated at its start $X$ and at the
local cap $L_{\rm loc}=585$ (\cref{lem:record:local-domain}).

\subsection{Middle blocks with $\sum\log p\le570$ and their degree sums}
\label{subsec:record:middle_blocks}

This subsection fixes the partition of the middle primes $p\equiv3\pmod4$
(called class-three primes below; $p\equiv1\pmod4$ are class-one) into
blocks and bounds the three sums
\[
 b_d(X,Y)=\sum_{\Blk\subseteq[X,Y)}q_{\Blk}^{2a}\,b_{\Blk}^{-d/2}
 \qquad(d=2,3,4)
\]
over the blocks $\Blk$ of a middle piece $[X,Y)$. Multiplied by the local
coefficients evaluated at $X$, these sums give every middle price
(\cref{lem:record:middle-prices}). The subsection also gives the lower
bounds for $\sum_{\Blk}b_{\Blk}^{-1}$ used by the mean credit
(\cref{lem:record:lower-start-sums}) and the caps of the selection factors of
the pieces used by the restorations (\cref{lem:piece-caps}).

Below the explicit Chebyshev threshold the partition consists of exact
blocks built from an authenticated prime histogram; above it, of
geometrically spaced cells partitioned into blocks of fixed cardinality $n_j$. Every block has logarithmic length at most the
budget $L=570$, so its selection factor is at most $e^{2aL}$; the local laws
of a piece are evaluated at the larger local cap $L_{\rm loc}=585$. The
partition depends on no parameter of the induction and is fixed before the
scale $N$.

Throughout, $\theta(t;4,3)=\sum_{p\le t,\ p\equiv3\ (4)}\log p$, and
$\mathsf e_X=2411/(5000000\log X)$ is the relative error of \cref{thm:bmor}.
We put $X_{\rm Ch}=4\,800\,162\,890$; thus $X_{\rm Ch}-1$ is the threshold of
that theorem, and $|\theta(t;4,3)-t/2|\le\mathsf e_Xt$ for all
$t\ge X\ge X_{\rm Ch}$. The threshold has no slack: an analytic cell may not
start below $X_{\rm Ch}$.

\begin{definition}[Exact log-budget blocks]\label{def:record:exact-blocks}
Let $L=570$ and let
\[
 \mathfrak u_*=2^{40}L=626\,721\,627\,832\,320
\]
be the integer budget.
\begin{enumerate}[label=(\alph*)]
\item \emph{Pieces.} The middle range
 $[Q_3,P_3)=[19\,951\,586,\,167\cdot2^{40})$ is the union of the $33$ pieces
 $[X_i,Y_i)=[X_i,X_{i+1})$, $1\le i\le33$, whose endpoints
 $X_1<X_2<\dots<X_{34}$ are
 \begin{gather*}
  19951586,\ 21757358,\ 23726568,\ 25429506,\ 25874006,\ 28215802,\\
  2^{25},\ 39903170,\ 47453134,\ 56431604,\ 2^{26},\ 79806340,\
  94906266,\\ 112863208,\ 2^{27},\ 159612678,\ 189812532,\ 225726414,\
  2^{28},\\ 319225356,\ 379625064,\ 451452826,\ 2^{29},\ 638450710,\
  759250126,\\ 902905652,\ 2^{30},\ 2^{31},\ 2^{32},\ 2^{34},\ 2^{36},\
  2^{40},\ 2^{45},\ 167\cdot2^{40};
 \end{gather*}
 thus $X_1=Q_3$, $X_{33}=2^{45}$ and $Y_{33}=X_{34}=P_3$, the numbering of
 \cref{tab:record:middle-pieces}. A block of a piece contains only primes
 $p\equiv3\pmod4$ of that piece.
\item \emph{Mandatory threshold points.} For $97\le j\le128$ let
 $c_j=2\lceil2^{j/4}/2\rceil$ be the least even integer $\ge2^{j/4}$; let
 $c'_1=21\,757\,358$, $c'_2=25\,429\,506$ and $c'_3=25\,874\,006$, the least
 even integers $\ge2^{195/8}$, $2^{123/5}$ and $2^{197/8}$; and put
 $c_{129}=X_{\rm Ch}$. These $36$ numbers are the \emph{threshold points} (called cuts in the data; they are not mirror cuts), and the $35$ intervals between consecutive threshold points are the \emph{parts}. Every endpoint of (a) below $2^{32}$ is a threshold point ($Q_3=c_{97}$), so each piece below $2^{32}$ is a
 union of parts, and the piece $[2^{32},2^{34})$ contains the part
 $[c_{128},c_{129})$.
\item \emph{Exact blocks below $X_{\rm Ch}$.} Every prime
 $p\in[2^{24},X_{\rm Ch})$ carries an integer upper logarithm
 $\mathfrak u(p)$: the interval $[2^{24},X_{\rm Ch})$ is covered by $5668$
 fixed consecutive count bins $[{\rm lo},{\rm hi})$, and $\mathfrak u(p)$ is
 the least integer $\ge2^{40}\log{\rm hi}$ for the bin containing $p$, so
 that $\log p<\mathfrak u(p)/2^{40}$ (the bins and the authenticated
 histogram are in \cref{subsec:data:pieces}). In each part list its primes
 $p\equiv3\pmod4$ in increasing order. A prime joins the current block if
 the sum of $\mathfrak u$ over the block stays at most $\mathfrak u_*$;
 otherwise it starts a new block. The last block of a part is closed at the
 end of the part.
\item \emph{Fixed-cardinality cells above $X_{\rm Ch}$.} In a piece $[X,Y)$
 with $Y>X_{\rm Ch}$ put $x=\max(X,X_{\rm Ch})$,
 $\mathsf n=\lceil128\log_2(Y/x)\rceil$ and
 \[
  x_j=x^{1-j/\mathsf n}Y^{j/\mathsf n},\qquad
  n_j=\Bigl\lfloor\frac L{\log x_{j+1}}\Bigr\rfloor\qquad(0\le j<\mathsf n).
 \]
 In the cell $[x_j,x_{j+1})$ group its primes $p\equiv3\pmod4$, in
 increasing order, into consecutive blocks of $n_j$ primes; the last
 block of the cell may have fewer. The five pieces above $X_{\rm Ch}$ have
 $\mathsf n=236,256,512,640,306$; their integers $n_j$ are recorded in
 \cref{cert:record:middle-price}.
\end{enumerate}
The \emph{middle blocks} are the blocks of (c) and (d). Each lies in one
part or one cell, and therefore in one piece.
\end{definition}

No endpoint used in (c) or (d) is a prime: $X_{\rm Ch}$, the threshold points and every
$Y$ are even, and an interior $x_j$ cannot be an odd prime $p$, because
$p^{\mathsf n}=x^{\mathsf n-j}Y^j$ would be even. The cells and the
cardinalities are deterministic functions of $L$ and of the endpoints; the
bin integers are fixed rational upper logarithms. No numerical choice
depends on $a$ or on $N$.

\begin{remark}[Restarts and fixed cardinality]
\label{rem:middle-restarts}
A block never straddles a part or a cell: the partition restarts at every
cut. This costs one restart term per cell in \cref{lem:record:block-charge},
but it removes every analysis of blocks that cross a boundary between
pieces with different local coefficients, and it lets each piece carry its
own local law. Inside a cell all blocks except the last have exactly $n_j$
primes, so the start of every later block is controlled by the full block
that precedes it; this turns the sum over block starts into a sum over
primes, to which the explicit Chebyshev bound is applied once, on the whole
interval.
The pieces need not be equally spaced: the threshold points $c'_1,c'_2,c'_3$, together with
$c_{98}$, are endpoints $X_i$, so $[c_{97},c_{99})$ is the union of five
consecutive pieces, each a union of parts. Each piece carries its own prime
counts, product
budget and restart, endpoint and jump terms; the pieces are disjoint and cover
the range, so adding their charges drops no boundary term.
\end{remark}

\begin{lemma}[Domain of the local laws]\label{lem:record:local-domain}
Let $\Blk$ be a middle block in the piece $[X,Y)$. Then:
\begin{enumerate}[label=(\roman*)]
\item every prime of $\Blk$ satisfies $p\equiv3\pmod4$ and
 $X\le b_{\Blk}\le p<Y$;
\item $\log q_{\Blk}\le L$; hence
 $|\Blk|\le\lfloor L/\log X\rfloor\le n_X:=\lfloor L_{\rm loc}/\log X\rfloor$
 with $L_{\rm loc}=585$, $\sum_{p\in\Blk}1/p\le L/(X\log X)$ and
 $q_{\Blk}^{2a}\le e^{2aL}$;
\item each bound listed below, computed with threshold $Q=X$ and
 logarithmic block length $L_{\rm loc}$, holds for every prime and block as
 stated in its cited result, hence for $\Blk$; these are the
 caps of the cuts of the root and of the twelve-label graphs with their block
 inflation (\cref{lem:layer12-caps}; \cref{imp:local:collision-gap} on the pieces
 below $2^{30}$; \cref{cert:positive-root-caps} on the pieces from $2^{30}$), the cap $k_{\rm D}$ of the reference fold record
 (\cref{lem:old-f-law}), the row constants $f_j,h_j$, the completion weights,
 the one-stage coefficients $k_{\rm new}$ and $s_{\rm fold}$ and the
 coefficients of the middle prices (\cref{lem:front12:reference},
 \cref{lem:record:one-stage-folded-e}, \cref{lem:record:middle-prices} and
 \cref{lem:middle-comparison-prices}).
\end{enumerate}
\end{lemma}

\begin{proof}
(i) holds by construction. For (ii), a block of
\cref{def:record:exact-blocks}(c) satisfies
\[
 \log q_{\Blk}<\sum_{p\in\Blk}\mathfrak u(p)/2^{40}\le\mathfrak u_*/2^{40}=L .
\]
A block of (d) in the cell $[x_j,x_{j+1})$ has at most $n_j$ primes, all
below $x_{j+1}$, so $\log q_{\Blk}<n_j\log x_{j+1}\le L$ by the definition of
$n_j$. Every prime of $\Blk$ is at least $X$, so
$|\Blk|\log X\le\log q_{\Blk}\le L\le L_{\rm loc}$,
$\sum_{p\in\Blk}1/p\le|\Blk|/X$, and
$q_{\Blk}^{2a}=e^{2a\log q_{\Blk}}\le e^{2aL}$.

(iii) The row constants, $k_{\rm new}$, $s_{\rm fold}$ and the coefficients of
\cref{lem:record:middle-prices} are admissible functions in the sense of
\cref{lem:majorant}, of $r=p^{-1/2}$ and of the block count, with their powers
of $r$ cancelled (for the row constants see the proof of
\cref{lem:front12:reference}(i), for $k_{\rm new}$ and $s_{\rm fold}$ their
formulas). By \cref{lem:majorant}(b) the value of such a coefficient at $X$,
with the count $n_X\ge|\Blk|$ of (ii), bounds its value at $b_{\Blk}\ge X$. The
coefficients of \cref{lem:middle-comparison-prices} combine the weights and
$f_1,h_1$ with the constants $s_F$, $a_F$ and $\ell_{\tlaw}$ of the signed fold
records and the coefficients $s_E$, $a_E$ and $g_\Sigma(X)$, which are valid at
every block of the piece by \eqref{eq:signed-constants} with
\cref{prop:signed-bounds}, by \cref{lem:record:exact-marginal-e}(ii) and by
\cref{cor:merge-repair}. The caps
are constants valid for every prime at least the threshold at which they are
computed, and this threshold is the piece start $X$ itself: every piece start
is a threshold of \cref{cert:layer12-grid} and an endpoint of
\cref{cert:l20-grid}, and the caps of the positive roots are evaluated at $X$
(\cref{cert:positive-root-caps}). The block caps hold for every block of the
piece by (i) and (ii), since $L\le L_{\rm loc}$; the grid domain $Q\ge2^{24}$,
length at most $10000$, contains $(X,L_{\rm loc})$. The cap $k_{\rm D}$ holds
at every block of the piece by \cref{lem:old-f-law}. The completion weights are
fixed numbers whose inequalities involve only these inputs. That the local data
of each of the $33$ pieces were computed at $(X,L_{\rm loc})$ is part of
\cref{cert:record:middle-price}.
\end{proof}

The selection exponent $2aL$ exceeds one: it is
$1236349836/1220703125<1.013$ at the baseline exponent $a$, and
$253204431/250000000$ at the reference exponent $a_0=\PushAnchor$ of the tabulated
middle prices. No statement of the paper requires $2aL\le1$: the restorations
and variance credits take an arbitrary upper bound for the block selection
factor, and every block is priced at its own $q_{\Blk}^{2a}$. A middle price at
$a\ge a_0$ is at most $t=e^{2(a-a_0)L}$ times its value at $a_0$
(\cref{lem:record:repricing}), and the selection factor of one block is at most
$C=e^{2aL}$, which is below $2.753348440327$ at the baseline exponent (\cref{def:closing-quantities},
\cref{tab:closing-values}). The restorations use the sharper caps of the
individual pieces (\cref{lem:piece-caps}).

\begin{lemma}[Block charges]\label{lem:record:block-charge}
Let $a>0$ and $s>0$.
\begin{enumerate}[label=(\roman*)]
\item \emph{Cells.} Let $[x,Y)$ be an interval with $x\ge X_{\rm Ch}$ and
 $x,Y$ even, $x_j=x^{1-j/\mathsf n}Y^{j/\mathsf n}$ for $0\le j\le\mathsf n$,
 and let $n_0,\dots,n_{\mathsf n-1}\ge1$ be integers. Partition the primes
 $p\equiv3\pmod4$ of each cell $[x_j,x_{j+1})$ into consecutive blocks of
 $n_j$ primes, the last possibly shorter. Put
 \[
  \Pi_j=\exp(2an_j\log x_{j+1}),\qquad g_j=\Pi_j/n_j,\qquad
  I_j(s)=\int_{x_j}^{x_{j+1}}\frac{t^{-s}}{\log t}\,dt
 \]
 (numbers attached to the cells). Then, with
$\mathsf e=\mathsf e_x:=2411/(5000000\log x)$,
 \begin{equation}\label{eq:record:cell-charge}
 \begin{split}
  \sum_{\Blk}q_{\Blk}^{2a}b_{\Blk}^{-s}\le{}&\sum_{j<\mathsf n}\Pi_jx_j^{-s}
   +\Bigl(\frac12+\mathsf e\Bigr)\sum_{j<\mathsf n}g_jI_j(s)\\
   &+2\mathsf e\,\frac{g_0x^{1-s}}{\log x}
   +2\mathsf e\sum_{0<j<\mathsf n}(g_j-g_{j-1})_+\frac{x_j^{1-s}}{\log x_j}.
 \end{split}
 \end{equation}
\item \emph{Exact blocks.} For a block $\Blk$ of
 \cref{def:record:exact-blocks}(c) put
 $\mathfrak u(\Blk)=\sum_{p\in\Blk}\mathfrak u(p)$,
 $k_{\Blk}=\lceil2^{16}\mathfrak u(\Blk)/\mathfrak u_*\rceil$ and
 \[
  w_2(\Blk)=\Bigl\lceil\frac{2^{95}}{b_{\Blk}}\Bigr\rceil,\quad
  w_3(\Blk)=\Bigl\lceil\frac{2^{127}}{b_{\Blk}\lfloor(2^{64}b_{\Blk})^{1/2}
  \rfloor}\Bigr\rceil,\quad
  w_4(\Blk)=\Bigl\lceil\frac{2^{95}}{b_{\Blk}^2}\Bigr\rceil.
 \]
 Then $1\le k_{\Blk}\le2^{16}$,
 $q_{\Blk}^{2a}\le\exp(2a\mathfrak u_*k_{\Blk}/2^{56})\le e^{2aL}$ and
 $b_{\Blk}^{-d/2}\le w_d(\Blk)/2^{95}$. Consequently, for every part
 $\mathfrak p$,
 \[
  \sum_{\Blk\subseteq\mathfrak p}q_{\Blk}^{2a}b_{\Blk}^{-d/2}\le
  \sum_{k=1}^{2^{16}}\exp\Bigl(\frac{2a\mathfrak u_*k}{2^{56}}\Bigr)
  \frac{w_d(\mathfrak p,k)}{2^{95}},\qquad
  w_d(\mathfrak p,k)=\sum_{\Blk\subseteq\mathfrak p,\ k_{\Blk}=k}w_d(\Blk).
 \]
\end{enumerate}
In particular, for every middle piece $[X,Y)$ and $d\in\{2,3,4\}$,
$b_d(X,Y)$ is at most the sum of (ii) over the parts of the piece below
$X_{\rm Ch}$ and of \eqref{eq:record:cell-charge} with $s=d/2$ over its cells
above $X_{\rm Ch}$.
\end{lemma}

\begin{proof}
(i) \emph{Block starts.} Fix a cell $j$. Its first block starts at a prime
$b\ge x_j$, so it contributes at most $\Pi_jx_j^{-s}$: its primes are
below $x_{j+1}$, and it has at most $n_j$ of them, so
$q_{\Blk}<x_{j+1}^{n_j}$. Every later block $\Blk'$ of the cell is preceded by
a block $\Blk$ of exactly $n_j$ primes, all smaller than $b_{\Blk'}$. Since
$t^{-s}$ decreases, $b_{\Blk'}^{-s}\le n_j^{-1}\sum_{p\in\Blk}p^{-s}$.
Distinct blocks $\Blk'$ give distinct blocks $\Blk$, so
\[
 \sum_{\Blk\subseteq\text{cell }j}q_{\Blk}^{2a}b_{\Blk}^{-s}\le\Pi_jx_j^{-s}
  +g_j\sum_{\substack{x_j\le p<x_{j+1}\\ p\equiv3\ (4)}}p^{-s}.
\]
Empty cells and short last blocks are covered, and every cell keeps its
own restart term $\Pi_jx_j^{-s}$.

\emph{Partial summation.} Put $f_s(t)=t^{-s}/\log t$ and let
$\phi(t)=g_jf_s(t)$ for $x_j\le t<x_{j+1}$. Then $\phi$ is right-continuous
and positive, it decreases on each open cell (as $s>0$ and $t>1$), and its
jump at an interior endpoint is $\Delta_j=(g_j-g_{j-1})f_s(x_j)$. With
$\eta(t)=\theta(t;4,3)-t/2$,
\[
 \sum_jg_j\sum_{\substack{x_j\le p<x_{j+1}\\ p\equiv3\ (4)}}p^{-s}
 =\int_{[x,Y)}\phi\,d\theta(\cdot;4,3)
 =\frac12\int_x^Y\phi(t)\,dt+\int_{[x,Y)}\phi\,d\eta .
\]
No endpoint is a prime, so $\eta$ has no atom at any $x_j$ or at $Y$, and
Stieltjes integration by parts with the jumps of $\phi$ gives
\[
 \int_{[x,Y)}\phi\,d\eta=\eta(Y)\phi(Y^-)-\eta(x)\phi(x)-\int_x^Y\eta(t)\phi'(t)\,dt
 -\sum_{0<j<\mathsf n}\eta(x_j)\Delta_j ,
\]
the integral being over the open cells. By \cref{thm:bmor},
$|\eta(t)|\le\mathsf et$ on $[x,Y]$. Since $\phi\ge0$ and $\phi'<0$ on the
open cells, the modulus of the right side is at most
\[
 \mathsf e\Bigl[Y\phi(Y^-)+x\phi(x)+\int_x^Yt(-\phi'(t))\,dt
 +\sum_jx_j|\Delta_j|\Bigr].
\]
Integrating by parts on each open cell and summing,
$\int_x^Yt(-\phi')\,dt=x\phi(x)-Y\phi(Y^-)+\int_x^Y\phi+\sum_jx_j\Delta_j$.
Since $\Delta_j+|\Delta_j|=2(\Delta_j)_+$, the error is at most
$\mathsf e[2x\phi(x)+\int_x^Y\phi+2\sum_jx_j(\Delta_j)_+]$. Adding the main
term $\frac12\int\phi=\frac12\sum_jg_jI_j(s)$ and the restart terms gives
\eqref{eq:record:cell-charge}. The argument does not use monotonicity of
$g_j$, and it applies a single Chebyshev error to the whole interval; no additional error term arises at interior cell
endpoints beyond the jump terms displayed in
\eqref{eq:record:cell-charge}.

(ii) The bin integers satisfy $\log p<\mathfrak u(p)/2^{40}$, so
$\log q_{\Blk}<\mathfrak u(\Blk)/2^{40}$. A single prime has
$\mathfrak u(p)\le\lceil2^{40}\log X_{\rm Ch}\rceil<2.46\cdot10^{13}<\mathfrak u_*$,
so a block started by one prime is within the budget, and the greedy rule
gives $1\le\mathfrak u(\Blk)\le\mathfrak u_*$, hence $1\le k_{\Blk}\le2^{16}$
and $\mathfrak u(\Blk)\le\mathfrak u_*k_{\Blk}/2^{16}$, which gives the price
bound and $\exp(2a\mathfrak u_*k_{\Blk}/2^{56})\le\exp(2a\mathfrak u_*/2^{40})
\le e^{2aL}$. For $d=2,4$ the weight bound is immediate. For $d=3$ put
$s_0=\lfloor(2^{64}b)^{1/2}\rfloor\le2^{32}b^{1/2}$; then
$2^{127}/(bs_0)\ge2^{95}b^{-3/2}$. Group the blocks of the part by $k_{\Blk}$.
Every block, including the first block after each cut, is an actual block
of the partition, so no restart term is needed.
\end{proof}

\begin{corollary}[One positive functional for all charge degrees]
\label{cor:positive-charge-functional}
For fixed cells and fixed $a$, the right side of
\eqref{eq:record:cell-charge} is $\Phi(t^{-s})$, where
\begin{align*}
 \Phi(f)={}&\sum_j\Pi_j f(x_j)
 +(\tfrac12+\mathsf e)\sum_jg_j\int_{x_j}^{x_{j+1}}
                                      \frac{f(t)}{\log t}\,dt\\
 &+2\mathsf e\frac{g_0x}{\log x}f(x)
 +2\mathsf e\sum_{j>0}(g_j-g_{j-1})_+
                              \frac{x_j}{\log x_j}f(x_j).
\end{align*}
This functional is positive and independent of the degree $s$.
Hence any pointwise polynomial inequality in $r=t^{-1/2}$ on
$[x,Y]$ may be integrated against $\Phi$, including an
inequality with signed coefficients. A common positive repricing
factor preserves this property.
\end{corollary}
\begin{proof}
Every displayed weight is nonnegative, and substitution of $f(t)=t^{-s}$
gives \eqref{eq:record:cell-charge}. Apply positivity to the difference
of the two sides of the pointwise inequality.
\end{proof}

The inequality is applied before rounding. Independently rounded
upper bounds for different monomials do not define the same positive
functional and cannot be substituted into its negative coefficients.
For the exact blocks the same principle applies to the actual common
block sum; their independently rounded degree histograms supply upper
bounds, not signed functional identities.


The same partial summation, with the error subtracted, gives lower bounds
for the unpriced start sums.

\begin{lemma}[Lower start sums]\label{lem:record:lower-start-sums}
\leavevmode
\begin{enumerate}[label=(\roman*)]
\item In the setting of \cref{lem:record:block-charge}(i),
 \[
  \sum_{\Blk}b_{\Blk}^{-1}\ \ge\ \Bigl(\frac12-\mathsf e\Bigr)
   \sum_{j<\mathsf n}\frac1{n_j}\log\frac{\log x_{j+1}}{\log x_j}
   -2\mathsf e\Bigl(\frac1{n_0\log x}
   +\sum_{0<j<\mathsf n}\frac{(1/n_j-1/n_{j-1})_+}{\log x_j}\Bigr).
 \]
\item For every part $\mathfrak p$ of \cref{def:record:exact-blocks}(c),
 \[
  \sum_{\Blk\subseteq\mathfrak p}b_{\Blk}^{-1}\ \ge\
  \frac{\sum_{\Blk\subseteq\mathfrak p}w_2(\Blk)-\#\{\Blk\subseteq\mathfrak p\}}
  {2^{95}} .
 \]
\end{enumerate}
\end{lemma}

\begin{proof}
(i) A block of cell $j$ has at most $n_j$ primes, all at least its first
prime $b$, so $\sum_{p\in\Blk}p^{-1}\le n_jb^{-1}$. Summing over the blocks,
$\sum_{\Blk}b_{\Blk}^{-1}\ge\sum_j n_j^{-1}\sum_{p\in\text{cell }j}p^{-1}$.
Apply the partial summation of the proof of
\cref{lem:record:block-charge}(i) with $s=1$ and $g_j$ replaced by $1/n_j$:
the main term is $\frac12\sum_jn_j^{-1}I_j(1)$ with
$I_j(1)=\log(\log x_{j+1}/\log x_j)$, and the same error bound, now
subtracted, gives (i). No first-block term is subtracted, and the bound
holds for arbitrary $n_j$. (ii) $w_2(\Blk)<2^{95}/b_{\Blk}+1$.
\end{proof}

The lower sum of (i) is a bound for a whole interval: it does not assume
that a single fine cell has relative prime density close to one half.

One budget $L$ bounds every block, but the largest block product of a piece
may be smaller; the restorations of \cref{sec:variance} use the following caps,
one for each piece.

\begin{lemma}[Piece caps]\label{lem:piece-caps}
For a middle piece $[X_i,Y_i)$ let $L_i$ be the largest of the numbers
$k_{\Blk}L/2^{16}$ over its blocks $\Blk$ of
\cref{def:record:exact-blocks}(c) and of the recorded upper bounds for
$n_j\log x_{j+1}$ over its cells of \cref{def:record:exact-blocks}(d). Then
every block of the piece has $\log q_{\Blk}\le L_i\le L$. Consequently, for
$0<a\le a_*$ and $C_i=e^{2a_*L_i}$, every block of piece $i$ has
$q_{\Blk}^{2a}\le C_i$, and distinct blocks $\Blk_1,\dots,\Blk_k$ lying in
pieces $i_1,\dots,i_k$ have $\prod_lq_{\Blk_l}^{2a}\le C_m:=\prod_lC_{i_l}$
($m$ the multiset of the pieces $i_1,\dots,i_k$).
\end{lemma}

\begin{proof}
For a block of (c), $\log q_{\Blk}<\mathfrak u(\Blk)/2^{40}\le
k_{\Blk}\mathfrak u_*/2^{56}=k_{\Blk}L/2^{16}$ by
\cref{lem:record:block-charge}(ii), as $\mathfrak u_*=2^{40}L$; and
$k_{\Blk}\le2^{16}$. A block of (d) in the cell $[x_j,x_{j+1})$ has at most
$n_j$ primes, all below $x_{j+1}$, so $\log q_{\Blk}<n_j\log x_{j+1}$, and
$n_j\log x_{j+1}\le L$ by the definition of $n_j$. The authenticated histogram
stores $k_{\Blk}$ for every exact block, so the maximum over the piece bounds
every block of it; it is not an average. The products follow from
$q_{\Blk}^{2a}\le e^{2aL_i}\le e^{2a_*L_i}$.
\end{proof}

The numbers $L_i$ lie between $556.79$ and $570$
(\cref{tab:data:piece-values}). They are recorded in
\path{research/push_2026_10_25/artifacts/author.json} (key
\texttt{restoration\_product\_caps}; producer
\path{research/advance_2026_10_24/envelopes/piece_product_caps.py}, which
authenticates the histogram and the partition run). The restorations
(\cref{subsec:variance:final}) use them at $a_*=a$.

\Cref{cert:record:middle-price} (\cref{app:certificates}) records the finite facts about
this partition (the tiling, the bins and the histogram of the exact blocks, the
cells and their integers $n_j$, the piece log caps $L_i$ and the lower start
sums of \cref{lem:record:lower-start-sums}) and, at the reference exponent
$a_0$, the four totals of \cref{subsec:middle-frame-prices}. The degree charges
of each piece are listed in \cref{tab:record:middle-pieces}, and its four sums
in \cref{tab:record:middle-charges}.

\begin{remark}[What rests on the explicit Chebyshev bound]
\label{rem:record:middle-scope}
In the middle range \cref{thm:bmor} enters only through
\eqref{eq:record:cell-charge} and \cref{lem:record:lower-start-sums}(i): it
prices the five analytic intervals above $X_{\rm Ch}$, which carry about
two thirds of the total degree-two charge at the reference exponent $a_0$
(pieces $29$--$33$ of \cref{tab:record:middle-pieces}), and it gives the
analytic lower start sums of the mean credit. Below $X_{\rm Ch}$ the blocks
are built from exact prime counts, and no Chebyshev estimate is used.
\end{remark}

\subsection{Ordered pairs at $p\equiv3\pmod4$}\label{subsec:middle-pairs}

Let $p\equiv3\pmod4$ be prime and $\chi$ its quadratic character, and write
$x\to y$ if $\chi(y-x)=1$. Since $\chi(-1)=-1$ this is a tournament on
$\Fp$ (the Paley tournament).

\begin{lemma}[Class-three middle pairs]\label{lem:record:class-three}
Let $p\equiv3\pmod4$ be prime.
\begin{enumerate}[label=(\roman*)]
\item For distinct $x,y\in\Fp$ there are exactly $(p-3)/4$ labels $z$ with
 $x\to z$ and $y\to z$, and exactly $(p-3)/4$ with $z\to x$ and $z\to y$.
\item Call $(x_0,x_0',x_2,x_2',u,u')$, six distinct labels, an
 \emph{$F$-configuration} if $y\to z$ for all $y\in\{x_0,x_0'\}$ and
 $z\in\{x_2,x_2'\}$. For $p\ge11$ every ordered pair $(u,u')$ of distinct
 labels is the middle pair of exactly $(p-3)(p-7)(p-4)(p-5)/16$
 $F$-configurations.
\item Let $p>15$. Fix distinct $u,u'$ and an ordered pair $(x_0,x_0')$ of
 distinct labels outside $\{u,u'\}$. The number of ordered pairs
 $(x_2,x_2')$ completing an $F$-configuration lies between
 $(p-11)(p-15)/16$ and $(p-3)(p-7)/16$. Consequently the conditional laws
 of $(x_0,x_0')$ and of $(x_2,x_2')$ given $(u,u')$ under
 $\lambda_p(G_{\rm F,2})$ (\cref{subsec:push:layer12spectral}), relative to the uniform
 distinct-pair law avoiding the middle pair, have densities in
 $[R_2^{-1},R_2]$ with $R_2=(p-3)(p-7)/((p-11)(p-15))$. For one layer the
 same holds with $R_1=(p-1)/(p-3)$.
\end{enumerate}
\end{lemma}

\begin{proof}
(i) is part of \cref{lem:paley-triple}.

(ii) By (i) there are $p(p-1)\lambda_0(\lambda_0-1)$ ordered outer
$K_{2,2}$ embeddings $(x_0,x_0',x_2,x_2')$, with $\lambda_0=(p-3)/4$;
each allows $(p-4)(p-5)$ ordered middle pairs. The maps
$z\mapsto\alpha z+\beta$ with $\chi(\alpha)=1$ preserve $F$-configurations
and act transitively on the ordered distinct pairs $(u,u')$ with fixed
$\chi(u'-u)$. The middle labels are unconstrained, so $(u,u')$ and $(u',u)$
have the same number of completions; since $\chi(-1)=-1$, this exchange
interchanges the two orbits. Hence the count is constant, equal to
$\lambda_0(\lambda_0-1)(p-4)(p-5)=(p-3)(p-7)(p-4)(p-5)/16$.

(iii) By (i), the common successors of $x_0,x_0'$ number $\lambda_0$;
removing those among $u,u'$ leaves $\lambda_0-h$ with $0\le h\le2$. Ordered
pairs of distinct ones number $(\lambda_0-h)(\lambda_0-h-1)$, which lies
between $(\lambda_0-2)(\lambda_0-3)$ and $\lambda_0(\lambda_0-1)$. The
conditional density of $(x_0,x_0')$ is its count divided by the mean count,
so it lies in $[R_2^{-1},R_2]$. The same argument with common predecessors
gives the opposite outer marginal. For one layer the successors of $x_0$
other than the middle label number $(p-1)/2$ or $(p-3)/2$, whence $R_1$.
\end{proof}

\begin{remark}[Why the middle range is class three]\label{rem:middle-class-three}
Part (ii) fails for $p\equiv1\pmod4$: at $p=13$ the two orbits of ordered
middle pairs have $312$ and $264$ completions. Part (ii) gives the $F$ and
$E$ laws the same middle law, so that in the completion and in the first
stage the $F$ and $E$ kernels are densities against the one conditional
reference of part~(iii) (the uniform distinct-pair law given the middle
pair) and the fold kernel has middle-pair density identically $1$
(\cref{rem:polarization-uses}); this is why the middle range consists of
class-three primes, every class-one prime below $P_4$ being a prefix prime
(\cref{lem:frontier:iidlocal}).
\end{remark}

\subsection{The twelve-label graphs and bounds $-gp^{-1/2}I\preceq A_p\preceq(1+c^+/p)I$}
\label{subsec:push:layer12spectral}

On a block $\Blk$ of class-three primes we write $x\to y$ if $\chi_p(y-x)=1$
at every prime $p$ of $\Blk$. The \emph{twelve-label root} $G_{\rm L,2}$, the
\emph{$F$ graph} $G_{\rm F,2}$ and the \emph{$E$ graph} $G_{\rm E,2}$ are the
graphs $G_{\rm L}$, $G_{\rm F}$, $G_{\rm E}$ of
\cref{lem:record:fold-identification} with the layer set $\{0,1\}$, and
$G_{\rm L,1}$, $G_{\rm F,1}$, $G_{\rm E,1}$ are their one-layer versions (layer
set $\{0\}$). Thus $G_{\rm L,2}$ has the vertices
$(u,\ell)\in\Z/6\Z\times\{0,1\}$ and the edges $(u,\ell)\to(\rot u,\ell')$ for
all $u,\ell,\ell'$, the four layer lifts of each of the six rotation edges
$0\to2\to4\to0$ and $1\to5\to3\to1$: $12$ vertices, $24$ edges, every degree
$4$. The $F$ and $E$ graphs have three \emph{objects} $0,1,2$, each carrying
one label per layer; the $F$ graphs have all the edges from the labels of
object $0$ to those of object $2$ (four, respectively one), the $E$ graphs no
edges. For an oriented graph $G$,
$\lambda_{\Blk}(G)=\bigotimes_{p\in\Blk}\lambda_p(G)$ is the product of the
uniform laws on the injective strict-square labellings
(\cref{def:graphs:labeling}).

\begin{lemma}[Twelve-label geometry]\label{lem:layer12-geometry}
\leavevmode
\begin{enumerate}[label=(\alph*)]
\item The chamber reflections $u\mapsto k-u$ ($k\in\{1,3,5\}$), the translation
 $u\mapsto u+2$ and the layer swap $\ell\mapsto1-\ell$ map the edge set of
 $G_{\rm L,2}$ onto itself; the first two do the same for $G_{\rm L,1}$; the
 layer swap is an automorphism of $G_{\rm F,2}$ and $G_{\rm E,2}$.
\item The subgraph of $G_{{\rm L},j}$ induced on the objects $0,1,2$ of the
 plus half $H_0$ is $G_{{\rm F},j}$, $j=1,2$.
\item The subgraph of $G_{\mathfrak f,2}$ induced on the layer $\ell=0$ is
 $G_{\mathfrak f,1}$, $\mathfrak f\in\{\rm L,F,E\}$.
\item $G_{{\rm F},j}$ and $G_{{\rm E},j}$ induce the same graph, without
 edges, on the object pairs $\{0,1\}$ and $\{1,2\}$.
\end{enumerate}
\end{lemma}

\begin{proof}
The chamber maps of (a) and part (b) are
\cref{lem:record:fold-identification}(a),(c); the layer swap permutes the
four layer lifts of every edge. (c) and (d) are read off from the
definitions. The script \path{research/push_2026_10_17/sparse/layer12.py}
reconstructs these facts.
\end{proof}

The spectral caps concern five cuts: the first-coordinate cut $\vartheta_0$
of $G_{\rm L,2}$ and of $G_{\rm L,1}$ (named L20 and L10), and the layer
cut $\iota$ of $G_{\rm L,2}$, $G_{\rm F,2}$ and $G_{\rm E,2}$ (named L21, F21
and E21). Each is spectral data in the sense of \cref{def:spec:data}, with the
deletion order $\mathfrak o$ of that definition fixed as in the data of
\cref{cert:layer12-grid}; by \cref{lem:layer12-geometry}(b),(c) its half
graph $G^+$ is $G_{{\rm F},j}$ for L20 and L10 and the one-layer graph for
the three layer cuts. For a prime $p$ let $\mathsf A_p$ be the normalized
cut operator of \cref{lem:spec:cutkernel} on $L^2(\lambda_p(G^+))$, so that
$\E_{\lambda_p(G)}[\Phi(x)\Phi(x')]=\ip\Phi{\mathsf A_p\Phi}$ for every real
function $\Phi$ on the half, $x$ and $x'$ being the half and the reflected
half of a labelling; on a block, $\mathsf A_{\Blk}=\bigotimes_{p\in\Blk}\mathsf A_p$,
and the same identity holds for $\lambda_{\Blk}(G)$
(\cref{lem:spec:cutkernel}(e)).

Route~A (\cref{thm:spectral-criterion}) gives the caps of L21, E21, L10 and F21
(\cref{lem:layer12-caps}). The cap of L20 of the uniform root is proved directly
from complete actual stars (\cref{subsec:middle-l20}); on the seven pieces from $2^{30}$ the caps of L20 and L21 of the positive root are those of
\cref{sec:positive-roots}. For the folds the reference uses the cap $k_{\rm D}$
of \cref{lem:old-f-law}, not the F21 cap.

\begin{lemma}[Twelve-label caps]\label{lem:layer12-caps}
Let $X$ be the start of a middle piece, and for one of the cuts L21, E21, L10
and F21 let $g_X$ be the class-three negative constant and $c^+_X$ the positive
constant recorded at the threshold $X$ in \cref{cert:layer12-grid}.
\begin{enumerate}[label=(\alph*)]
\item For every prime $p\ge X$ with $p\equiv3\pmod4$,
 $-g_Xp^{-1/2}I\preceq\mathsf A_p\preceq(1+c^+_X/p)I$.
\item Let $\Blk$ be a middle block, with least prime $b$, of the piece starting
 at $X$. Then
 $\mathsf A_{\Blk}\succeq-g_X\exp\bigl(c^+_XL_{\rm loc}/(X\log X)\bigr)\,b^{-1/2}$.
 The coefficient $g_X\exp(c^+_XL_{\rm loc}/(X\log X))$ is the \emph{block cap}
 of the cut for the piece.
\end{enumerate}
The reference uses the block caps of L21 on the pieces below $2^{30}$ and those
of E21 and L10 on every piece; the cap of F21 is recorded but not used.
\end{lemma}

\begin{proof}
(a) This is route~A, \cref{thm:spectral-criterion}(a),(d) for the class
$\mathfrak c=3$, with $p_0=X$ and $c^+_X$ the constant $B^+$ of (d). The
spectral data satisfy (S1) at $X$: the graphs have at most twelve vertices
(half-graphs on $n\le6$ vertices, in the notation of
\cref{thm:spectral-criterion}) and maximum degree at most four, and $X\ge Q_3>2^{24}$, so
$e_\Delta(X,2n)\le e_4(2^{24},12)<5\cdot10^{-3}$ and $X>4n^2$. The inputs
(R1)--(R3) are the standard ones of \cref{rem:spectral-criterion:uses}(i): the
nonstar bound of the compressed matching partition for these spectral
data (\cref{cor:extraction-compressed}), the natural coincidence bound
\eqref{eq:collision-natural} of \cref{lem:oct24-collision}, and the standard
star and coupling inputs. The evaluator verifies the test
\eqref{eq:spectral-test} in exact rational arithmetic at $X$ for both
class-three signs, with one common value of $\tau^{\rm sp}_{3,\pm}$, and
records $g_X$, rounded upward, as the class-three constant $g_3$ of the
theorem; it computes $c^+_X$ with the largest $g$ of all four tests, which
satisfies $gX^{-1/2}<1$ and is at least $g_3$, as (d) allows
(\cref{rem:spectral-criterion:certificate}).
(b) By \cref{lem:record:local-domain}(i),(ii), the primes of $\Blk$ satisfy
$p\equiv3\pmod4$ and $p\ge X$, and $\sum_{p\in\Blk}\log p\le L\le L_{\rm loc}$.
Apply \cref{cor:spectral-criterion-blocks} with its $P$ equal to $X$ and its
length $L$ equal to $L_{\rm loc}$.
\end{proof}

\begin{certificate}[The twelve-label grids]\label{cert:layer12-grid}
For each of the five cuts of the twelve-label law and each of the
$35$ thresholds listed below, the constants of the Route~A
criterion (\cref{thm:spectral-criterion}) have been evaluated and
its four exact positive-semidefiniteness tests
\eqref{eq:spectral-test} pass; every middle piece start is one of
the thresholds, so these give the caps at every piece. The values
are stored in two data files, built with the spectral evaluator
\path{research/push_2026_10_17/sparse/layer72_spectral.py} applied to the
twelve-label graphs (\path{research/push_2026_10_17/sparse/layer12_spectral.py}).
The record \path{research/frontier_2026_10_17/root/middle12_grid.json}, built by
\path{research/frontier_2026_10_17/root/build_middle12_grid.py}, contains $26$
thresholds: $2^{24}$, $2^{25}$, $2^{26}$, the sixteen piece starts from
$79\,806\,340$ to $2^{30}$, and $2^{31}$, $2^{32}$, $2^{34}$, $2^{36}$,
$2^{40}$, $2^{45}$, $2^{50}$; the replay record
\path{research/frontier_2026_10_17/root/middle12_grid_replay.json} repeats its
$520$ tests with fresh partitions. The record
\path{research/continuation_2026_10_24/global/middle_grid_new.json}, built by
\path{research/continuation_2026_10_24/global/middle_grid.py}, contains the nine
piece starts $19\,951\,586$, $21\,757\,358$, $23\,726\,568$, $25\,429\,506$,
$25\,874\,006$, $28\,215\,802$, $39\,903\,170$, $47\,453\,134$,
$56\,431\,604$; the review record
\path{research/continuation_2026_10_24/global/middle_grid_review.json} repeats
its $180$ tests with fresh partitions. The two records contain the same
compressed matching partition of each cut, which each review rebuilds, and,
for each threshold and cut, the constants
of \cref{thm:spectral-criterion} and four exact rational positive
semidefiniteness tests of \eqref{eq:spectral-test} (both real class-one signs
and both Hermitian class-three signs); in every record the two class-three
tests use the same $\tau^{\rm sp}$. Every piece start is a threshold of one of
the records: pieces $1$--$6$ and $8$--$10$ read the second, pieces $7$ and
$11$--$33$ the first. The final source stores for each piece and cut the
threshold $X$, the length $585$, the class-three $g$, $c^+$ and the block cap;
the native check
(\path{research/advance_2026_10_24/global/check_low_reference_native.py}) binds
every stored route-A spectrum to the record entry at $X$, and the L20 spectrum
to the endpoint $X$ of \cref{cert:l20-grid}, and recomputes the block caps; the
native check of the final source verifies that its spectra are those of the
source so checked, with L20 replaced by the protected collision cap
at every start and with both L20 and L21 replaced on the seven middle pieces
from $2^{30}$ by \cref{cert:positive-root-caps}. At $Q_3$ the
class-three values of $g$ are $2.57407111$ (L21), $0.00454092$ (E21) and
$1.75129044$ (L10), with block caps
below $2.574071293695$, $0.004540920001$ and $1.751290440133$. The records also
contain route-A constants of L20 and F21, which the reference does not use
(\cref{subsec:middle-l20}; for the folds \cref{lem:old-f-law},
\cref{cert:record:old-f-cap}), and of L21 at the starts from $2^{30}$, where the
root is a positive tilt (\cref{sec:positive-roots}). The finite
checks are inputs of the spectral proof, not substitutes for its graph
decomposition. The data are listed in \cref{tab:data:layer12-grid}.
\end{certificate}

\subsection{The chamber symmetry L20 for the uniform principal law}\label{subsec:middle-l20}

One cut governs the middle range: the chamber reflection in the
first coordinate, called L20 in the tables and certificates. The
lower spectral bound (the \emph{cap}) of its operator for the
twelve-label principal law is the most negative in the middle
range, so it is the bound the completion has to compensate. On the
middle intervals (``pieces'') starting below $2^{30}$ the
principal law is the uniform law $\lambda_{\Blk}(G_{\rm L,2})$; the
tilted laws of \cref{sec:positive-roots} start at $2^{30}$. This
subsection bounds the L20 operator for the uniform law directly,
by computing its star part exactly in its own half law
(\cref{subsec:spec:actual-stars}) instead of using the general
criteria of \cref{sec:spectral}. Fix
a prime $p\equiv3\pmod4$. By \cref{lem:layer12-geometry}(b) the half graph of
L20 is $G_{\rm F,2}$, so the half reference $\mu_p=\lambda_p(G_{\rm F,2})$ is
the uniform law on injective tuples $x=(a_0,a_1,b_0,b_1,c_0,c_1)$ with
$a_i\to c_j$ for all $i,j$ (the labels of the objects $0,1,2$ in the layers
$0,1$; coordinates numbered $0,\dots,5$ in this order). Let $T=\mathsf A_p$ be
the L20 cut operator on $L^2(\mu_p)$, $\mathcal I$ the subspace of functions
invariant under the simultaneous translations $x_k\mapsto x_k+y$
($y\in\Fp$), and $\chi^+_p=D_{\rm full}^2/L_{\rm low}^2$ the second bound of
\cref{lem:record:root-marginal}(ii), with $D_{\rm full},L_{\rm low}$ as defined
there. As $T\one$ is the density of the half law
of the root relative to $\mu_p=\lambda_{F,p}$ (\cref{lem:spec:cutkernel}(b)),
$\norm{T\one-\one}^2=\chi^2(\nu_{\mathrm L,p}\,\|\,\lambda_{F,p})\le\chi^+_p$,
where $\nu_{\mathrm L,p}$ is the $\vartheta_0$-half law of $\lambda_p(G_{{\rm L},2})$.

The cross pairs of the cut are the sixteen layer lifts of the four root edges
$b'\to a$, $c\to b'$, $b\to a'$ and $c'\to b$ between the half $x$ and the
reflected half $x'$. The \emph{soft L20 kernel} is the product over the
sixteen cross pairs of the factors $1+(D_{20})_{ij}\chi(x'_j-x_i)$, without the
indicator that $x$ and $x'$ have no common label (its factors are the signed
factors of the cut, not the soft factors of \cref{cor:soft-kernels}), with the
signed cross matrix
\begin{equation}\label{eq:l20-cross-matrix}
 D_{20}=\begin{pmatrix}
 0&0&-1&-1&0&0\\0&0&-1&-1&0&0\\
 1&1&0&0&-1&-1\\1&1&0&0&-1&-1\\
 0&0&1&1&0&0\\0&0&1&1&0&0
 \end{pmatrix}
\end{equation}
(rows indexed by the coordinates of $x$, columns by those of $x'$; an edge
$y\to y'$ means $\chi(y'-y)=1$, and $\chi(-1)=-1$). The matrix $D_{20}$ is
antisymmetric. In the notation of \cref{lem:spec:expansion} the cross signs are
$\varsigma_{ij}=(D_{20})_{ij}$, so $T=K_0/Z_0$ and $K_0=K_{\rm soft}-T_{\rm coll}$,
where $K_{\rm soft}=\one\otimes\one+\sum_{S\ \rm star}T_S+T_{\rm ns}$ is the soft
kernel as an operator on $L^2(\mu_p)$. We call $C=(a_0,a_1,c_0,c_1)$
(coordinates $0,1,4,5$) the \emph{core} of $x$ and $b=(b_0,b_1)$ (coordinates
$2,3$; not the least prime $b_{\Blk}$) its \emph{pair}, write $\mathcal H=L^2_0(\Fp)$ and $\Phat=P/r$ for the
normalized Paley operator, where $P$ is the Paley operator of
\cref{lem:square-shift}, and use the features $V_k$, $W_S$ of
\cref{subsec:spec:actual-stars} relative to $\mu_p$ and their synthesis
$\mathcal F=(V_k,(W_S/r)_S)$. The signs of the core coordinates are
$s=(s_0,s_1,s_4,s_5)=(-1,-1,1,1)$, as in \cref{lem:record:weighted-moment}.

\begin{lemma}[Complete stars of L20]\label{lem:l20-star-matrix}
Let $p\equiv3\pmod4$ and $p>15$.
\begin{enumerate}[label=\textup{(\alph*)}]
\item The leaf supports with at least two elements of the stars of
 $K_{\rm soft}$ are the twelve sets $01$, $014$, $0145$, $015$, $04$, $045$,
 $05$, $14$, $145$, $15$, $23$, $45$, the subsets with at least two elements of
 the row supports $\{2,3\}$ and $\{0,1,4,5\}$ of $D_{20}$. With the six
 singletons, the Hermitian matrix $M(\Phat)$ of \cref{cor:star-matrix}
 satisfies
 \[
  \det(\lambda I-M(i))=\lambda^{14}(\lambda^2-4)(\lambda^2-30);
 \]
 in particular $\norm{M(i)}=\sqrt{30}$.
\item Under $\mu_p$ the conditional probability of each value of a core
 coordinate given the other five coordinates is at most $D_j/p$ with
 $D_j=4/(1-15q)$, and that of a pair coordinate is $D_j/p$ with
 $D_j=1/(1-5q)$. Let $\Gamma_{20}$ be one plus the largest row
 sum of the symmetric $18\times18$ matrix, indexed by the six
 singletons $V_k$ and the twelve supports $S$ of part~(a) (the
 features $W_S/r$), whose diagonal entries are $0$ and whose
 off-diagonal entries are: for two supports $S\ne T$, the bound of
 \cref{lem:general-star}(c) with the pivot $j\in S\triangle T$ of
 least $D_j$; for a singleton $k$ and a support $S$, the bound of
 \cref{lem:general-star}(d); and for two singletons, the bound of
 \cref{lem:general-star}(c) applied to their pair.
 Then $\norm{\mathcal F}^2\le\Gamma_{20}$; $\Gamma_{20}$ is a nondecreasing
 function of $r$ (with $q=r^2$), and $\Gamma_{20}\to1$ as $p\to\infty$.
\item For a label $z$ put
 \begin{gather*}
  \mathsf K(z,C)=\prod_{i=0,1}\bigl(1-\chi(z-a_i)\bigr)
  \prod_{j=0,1}\bigl(1+\chi(z-c_j)\bigr),\\
  \ell(z,C)=-\chi(z-a_0)-\chi(z-a_1)+\chi(z-c_0)+\chi(z-c_1),
 \end{gather*}
 $\mathsf K_\star=\mathsf K-1-\ell$ and $(W_\star h)(x)=\E_z\mathsf K_\star(z,C)h(z)$
 for $h\in\mathcal H$, and put
 \begin{gather*}
  \mathbf u=\tfrac12(V_0+V_1-V_4-V_5),\qquad
  \mathbf v=\tfrac1{\sqrt2}(V_2+V_3),\qquad
  \mathbf w=\tfrac12(V_0+V_1+V_4+V_5),\\
  \mathbf e=W_{23}/r,\qquad
  \mathbf z=\tfrac1{\sqrt{15}}\bigl(W_\star/r-2\mathbf u\Phat\bigr).
 \end{gather*}
 Then $\mathsf K=16\cdot\one\{c_0,c_1\to z\to a_0,a_1\}$; in particular
 $\mathsf K$ vanishes at the core labels, $\mathsf K^2=16\mathsf K$,
 $\mathsf K\ell=4\mathsf K$ and $\E_z\ell^2=4$. The function $\mathsf K_\star$ is
 the signed sum $\sum_S\sigma_S\psi_S$ of the eleven higher stars on the core,
 $\sigma_S=\prod_{k\in S}s_k$, $\psi_S(z,x)=\prod_{k\in S}\chi(z-x_k)$, and the
 doubly centred complete star part of $K_{\rm soft}$ is
 \begin{equation}\label{eq:l20-rank4-operator}
  T_\star=r\bigl\{\sqrt{30}(\mathbf v\mathbf z^*+\mathbf z\mathbf v^*)
  +2(\mathbf w\mathbf e^*+\mathbf e\mathbf w^*)\bigr\}.
 \end{equation}
\end{enumerate}
\end{lemma}

\begin{proof}
(c) For $z$ outside the core each factor of $\mathsf K$ is $0$ or $2$, and all
four are $2$ exactly when $\chi(z-a_i)=-1$ and $\chi(z-c_j)=1$ for all $i,j$,
that is, when $z\to a_i$ and $c_j\to z$. At $z=a_i$ the factor
$1+\chi(a_i-c_j)$ vanishes, because $a_i\to c_j$ and $\chi(-1)=-1$; at $z=c_j$
the factor $1-\chi(c_j-a_i)$ vanishes. Hence $\mathsf K\in\{0,16\}$,
$\mathsf K^2=16\mathsf K$, and $\ell=4$ where $\mathsf K\ne0$. Since
$\E_z\chi(z-y)\chi(z-y')=\one_{y=y'}-q$ \eqref{eq:two-root} and
$\sum_ks_k=0$, $\E_z\ell^2=4(1-q)+4q=4$. Expanding the product,
$\mathsf K=\sum_S\sigma_S\psi_S$ over all $S\subseteq\{0,1,4,5\}$; the empty set
gives $1$ and the singletons give $\ell$, so $\mathsf K_\star$ is the sum over
the eleven sets with at least two elements.

(a) and \eqref{eq:l20-rank4-operator}. By \cref{cor:star-matrix}, $M(\Phat)$ has
the singleton block $D_{20}\otimes\Phat$, the entry
$(-1)^{|S|}\prod_{j\in S}(D_{20})_{kj}$ at (singleton $k$, support $S$) and
$\prod_{i\in S}(D_{20})_{ij}$ at (support $S$, singleton $j$). The leaf sets of
the left stars are the subsets of the row supports of $D_{20}$, those of the
right stars the subsets of its column supports, and these are the same because
$D_{20}$ is antisymmetric; this gives the twelve supports. In $\C^{18}$, with
the coordinate vectors $e_k$ ($0\le k\le5$) and $e_S$, put, at the phase
$\Phat=i$,
\begin{gather*}
 \hat{\mathbf u}=\tfrac12(e_0+e_1-e_4-e_5),\quad
 \hat{\mathbf v}=\tfrac1{\sqrt2}(e_2+e_3),\quad
 \hat{\mathbf w}=\tfrac12(e_0+e_1+e_4+e_5),\quad
 \hat{\mathbf e}=e_{23},\\
 \hat{\mathbf z}=\tfrac1{\sqrt{15}}\Bigl(\sum_S\sigma_Se_S-2i\,\hat{\mathbf u}\Bigr),
\end{gather*}
the sum over the eleven supports contained in $\{0,1,4,5\}$. The four vectors
$\hat{\mathbf v},\hat{\mathbf z},\hat{\mathbf w},\hat{\mathbf e}$ are orthonormal
($|\hat{\mathbf z}|^2=(11+4)/15$), and
\[
 M(i)=\sqrt{30}\bigl(\hat{\mathbf v}\hat{\mathbf z}^*+\hat{\mathbf z}\hat{\mathbf v}^*\bigr)
 +2\bigl(\hat{\mathbf w}\hat{\mathbf e}^*+\hat{\mathbf e}\hat{\mathbf w}^*\bigr),
\]
as a comparison of entries shows: on the right the singleton entries are $i$ at
$(b_l,a_k)$ and $(c_k,b_l)$ and $-i$ at $(b_l,c_k)$ and $(a_k,b_l)$, which are
those of $iD_{20}$; the entries at $(b_l,S)$ and $(S,b_l)$, $S\subseteq
\{0,1,4,5\}$, are $\sigma_S=(-1)^{|S|}\prod_{j\in S}(D_{20})_{2j}
=\prod_{i\in S}(D_{20})_{i2}$; the entries at $(k,23)$ and $(23,k)$, $k$ a core
coordinate, are $1=(D_{20})_{k2}(D_{20})_{k3}$; and all other entries vanish on
both sides, since a core row of $D_{20}$ has the support $\{2,3\}$ and a pair row
the support $\{0,1,4,5\}$. So $M(i)$ has the eigenvalues $\pm\sqrt{30}$ and
$\pm2$ on the span of the four vectors and $0$ on its orthogonal complement,
which is (a). The synthesis maps $e_k$ to $V_k$ and $e_S$ to $W_S/r$, the phase
acting as $\Phat$ on the label input; it maps $\hat{\mathbf v}$,
$\hat{\mathbf w}$, $\hat{\mathbf e}$ to $\mathbf v$, $\mathbf w$, $\mathbf e$, and
$\hat{\mathbf z}$ to $\mathbf z$, since $\sum_S\sigma_SW_S=W_\star$ by (c). The
identity $T_\star=r\mathcal FM(\Phat)\mathcal F^*$ of \cref{cor:star-matrix} is
therefore \eqref{eq:l20-rank4-operator}.

(b) The coordinate $a_0$ is a common predecessor of $c_0,c_1$ distinct from
$a_1,b_0,b_1$. By \cref{lem:record:class-three}(i) two distinct labels have
exactly $(p-3)/4$ common predecessors, so given the other coordinates $a_0$ is
uniform on at least $(p-3)/4-3=(p-15)/4$ labels; the same holds for $a_1$ and,
with common successors of $a_0,a_1$, for $c_0,c_1$. A pair coordinate is
uniform on the $p-5$ labels outside the other five. This gives the $D_j$. The
off-diagonal Grams of $\mathcal F$ are bounded by the entries of the matrix
(\cref{lem:general-star}(c),(d)), the diagonal ones are at most one
(\cref{lem:general-star}(a)), and \cref{lem:general-star}(f) gives
$\norm{\mathcal F}^2\le\Gamma_{20}$. Every entry is built from $q$, $r$ and
$D_j-1\in\{(3+15q)/(1-15q),\,5q/(1-5q)\}$ by sums, products, square roots and
division by $1-q$, so it is admissible (\cref{lem:majorant}) and tends to $0$.
\end{proof}

The synthesis bound $\norm{\mathcal F}^2\le\Gamma_{20}$ alone
gives the coefficient $\sqrt{30}\,\Gamma_{20}$ of $r$ in the
negative cap; its excess over $\sqrt{30}$ is of order $r$, with a
large coefficient coming from the overlaps of the eighteen
features. The following propositions separate the
translation-invariant part of the kernel
(\cref{prop:l20-completion}), replace this synthesis bound by
the exact Gram matrices of the four
features of \eqref{eq:l20-rank4-operator}, so that the leading coefficient
becomes $\sqrt{30}+O(q)$ (\cref{prop:l20-rank4}), and factor the root through a
rectangular kernel whose nonstar part is priced quadratically
(\cref{prop:l20-rectangular,prop:l20-pair-gap}). What remains of order $r$ is the
collision cap (\cref{prop:l20-collisions} and the remark after it).

\begin{proposition}[Translation sectors]\label{prop:l20-completion}
Let $p\equiv3\pmod4$, $p\ge2^{20}$, let $N_{\rm ns}$ and $D_{\rm col}$ be raw nonstar
and collision caps of the soft L20 kernel (``raw'': bounds for the
unnormalized kernel $K_0$, before division by $Z_0$), $\norm{T_{\rm ns}}\le N_{\rm ns}$
and $\norm{T_{\rm coll}}\le D_{\rm col}$
(\cref{cor:extraction-compressed,lem:oct24-collision}), let $m_\star$ be the
averaged-star mean bound of \cref{lem:l20-star-mean}, and put
$\eps_Z=N_{\rm ns}+D_{\rm col}+m_\star$ and
$H_{20}=(N_{\rm ns}+D_{\rm col})/(1-\eps_Z)$, assuming $\eps_Z<1$. Then $T$
preserves $\mathcal I$ and $\mathcal I^\perp$, the ranges of the features
$V_k$, $W_S$ lie in $\mathcal I^\perp$, and the centred block of $T$ on
$\mathcal I\cap\one^\perp$ has norm at most $H_{20}$. If moreover $H_{20}<1$
and $C_\perp<1$ satisfy $T\succeq-C_\perp I$ and
$\norm{T|_{\mathcal I^\perp}}\le C_\perp$, then
\begin{equation}\label{eq:l20-caps}
 -g\,rI\preceq T\preceq(1+c^+/p)I,\qquad g=C_\perp/r,\qquad
 c^+=\frac{p\chi^+_p}{1-H_{20}}.
\end{equation}
\end{proposition}

\begin{proof}
By \cref{lem:spec:expansion}, $T=K_0/Z_0$ with $Z_0=\ip\one{K_0\one}_{\mu_p}$,
and by \cref{lem:l20-star-mean} the star means of $K_{\rm soft}$ sum to at most
$m_\star$ in modulus, so $Z_0\ge1-\eps_Z>0$. The simultaneous translations
preserve $\mu_p$ and $K_0$, so $T$ commutes with them and preserves
$\mathcal I$, and $\mathcal I^\perp$ as it is self-adjoint. The ranges of the
features lie in $\mathcal I^\perp$ by \cref{lem:general-star}. Let
$h,h'\in\mathcal I\cap\one^\perp$. Then $\E_{\mu_p}[h\mid x_k]=0$ for every
coordinate $k$ (\cref{lem:spec:affpair}(d)), so $T_Sh=0$ for every column star
$S$ (\cref{lem:spec:stars}(a)) and
$\ip{h'}{T_Sh}=\ip{T_{S^{\mathsf t}}h'}h=0$ for every row star
(\cref{lem:spec:stars}(b)); also $(\one\otimes\one)h=0$. By
\eqref{eq:record:root-kernel},
$\ip{h'}{K_0h}=\ip{h'}{(T_{\rm ns}-T_{\rm coll})h}$, of modulus at most
$(N_{\rm ns}+D_{\rm col})\norm h\norm{h'}$, and division by $Z_0$ bounds the
centred block by $H_{20}$. The vector $T\one-\one$ is translation invariant
and orthogonal to $\one$, as $\ip\one{T\one}=1$, and its squared norm is at
most $\chi^+_p$. Now \cref{prop:matching:positive}(a$''$), for the group of
translations, gives
$T\preceq\max\bigl(C_\perp,1+\chi^+_p/(1-H_{20})\bigr)I=(1+c^+/p)I$, as
$C_\perp<1$. The lower bound is the hypothesis.
\end{proof}

The lower bound is required on the whole space;
\cref{prop:l20-rectangular} supplies it.

\begin{proposition}[Gram bounds for the four L20 star features]\label{prop:l20-rank4}
Let $p\equiv3\pmod4$, $p\ge2^{20}$, let
$\bar{\mathsf K}(C)=\E_z\mathsf K(z,C)$, $m_{\mathsf K}=\E_{\mu_p}\bar{\mathsf K}-1$
and $v_{\mathsf K}=\E_{\mu_p}(\bar{\mathsf K}-1)^2$, let $m_\star$ be the bound
of \cref{lem:l20-star-mean}, and put
\begin{gather*}
 \mathfrak m_+=2q+\tfrac14m_\star,\qquad
 g_{\mathbf z}=\frac{1+\frac{14}{15}\mathfrak m_+}{1-q},\qquad
 \rho=\frac{(12+8r)q^2}{(1-4q)(1-5q)},\\
 \rho_{\mathbf v\mathbf z}=\frac{4\sqrt{2/15}\,qr(1+\mathfrak m_+)}{(1-4q)(1-q)},\qquad
 \rho_{\mathbf v\mathbf w}=\frac{2\sqrt2\,q}{1-q},\qquad
 \rho_{\mathbf v\mathbf e}=\frac{\sqrt2\,qr}{1-q},\\
 \rho_{\mathbf z\mathbf w}=\frac{2r(1+\mathfrak m_+)}{\sqrt{15}(1-q)},\qquad
 \rho_{\mathbf z\mathbf e}=\sqrt{g_{\mathbf z}}\,\rho,\qquad
 k_{\rm top}=\sqrt{30}\,(g_{\mathbf z}+\rho_{\mathbf v\mathbf z}),\qquad
 k_{\rm low}=2(1+\rho),\\
 k_\times^2=2\sqrt{30}\bigl(\rho_{\mathbf v\mathbf w}^2+\rho_{\mathbf v\mathbf e}^2
 +\rho_{\mathbf z\mathbf w}^2+\rho_{\mathbf z\mathbf e}^2\bigr),\qquad
 k_4=k_{\rm top}+\frac{k_\times^2}{\sqrt{30}-k_{\rm low}}.
\end{gather*}
Then $\sqrt{30}-k_{\rm low}>0$, $|m_{\mathsf K}|\le\mathfrak m_+$, and on
$\mathcal H$:
\begin{enumerate}[label=\textup{(\alph*)}]
\item $\mathbf u^*\mathbf u=\frac1{1-q}I$,
 $\mathbf w^*\mathbf w=\frac{1-4q}{1-q}I$,
 $\mathbf u^*\mathbf w=\frac{2r}{1-q}\Phat$ and
 $\mathbf v^*\mathbf v=\frac{1-2q}{1-q}I$;
\item $W_\star^*W_\star=\frac{q(11+6m_{\mathsf K}-v_{\mathsf K})}{1-q}I$,
 $\mathbf u^*W_\star=-\frac{2m_{\mathsf K}r}{1-q}\Phat$ and
 $\mathbf w^*W_\star=\frac{2q(1-m_{\mathsf K})}{1-q}I$;
\item $\mathbf z^*\mathbf z=\frac{15+14m_{\mathsf K}-v_{\mathsf K}}{15(1-q)}I
 \preceq g_{\mathbf z}I$,
 $\mathbf w^*\mathbf z=-\frac{2r(1+m_{\mathsf K})}{\sqrt{15}(1-q)}I$,
 $\mathbf v^*\mathbf z=\frac{4\sqrt{2/15}\,qr(1+m_{\mathsf K})}{(1-4q)(1-q)}I$ and
 $\mathbf v^*\mathbf w=-\frac{2\sqrt2\,q}{1-q}I$;
\item $\mathbf e^*\mathbf e=\frac{1-2q-q^2}{1-q}I$,
 $\mathbf v^*\mathbf e=\frac{\sqrt2\,qr}{1-q}I$,
 $\norm{\mathbf w^*\mathbf e}\le\rho$ and
 $\norm{\mathbf z^*\mathbf e}\le\rho_{\mathbf z\mathbf e}$.
\end{enumerate}
Consequently, if $\mathfrak F=(\mathbf v,\mathbf z,\mathbf w,\mathbf e)$ and
$\Xi$ is a real symmetric $4\times4$ matrix with
$-\Delta_\star\preceq\Xi\preceq\Delta_\star$,
$\Delta_\star=\operatorname{diag}(\sqrt{30},\sqrt{30},2,2)$, then
$\norm{\mathfrak F\Xi\mathfrak F^*}\le k_4$. In particular
$\norm{T_\star}\le k_4r$.
\end{proposition}

\begin{proof}
For $p\ge2^{20}$, $\rho<1/2$, so $k_{\rm low}<3<\sqrt{30}$.

\emph{The mean.} By \cref{lem:l20-star-matrix}(c),
$\bar{\mathsf K}-1=\sum_S\sigma_Sw_S$ over the eleven higher supports $S$ on the
core, with $w_S(C)=\E_z\psi_S(z,x)$. The six two-element supports have
$w_S=-q$ \eqref{eq:two-root} and $\sum\sigma_S=\frac12((\sum_ks_k)^2-4)=-2$,
which contributes $2q$. Under $\mu_p$ the core is uniform on the admissible
quadruples of \cref{lem:record:weighted-moment} (each has exactly $(p-4)(p-5)$
completions by the pair), whose probability under independent uniform labels
is $Z_4/16$ with $Z_4=(1-q)(1-3q)(1-7q)$, while $P_{\rm adm}D_4=16$ on
them ($D_4$ the indicator that the four labels are distinct, and $P_{\rm adm}$
the product of \cref{lem:record:weighted-moment}, $16$ on admissible
quadruples and $0$ on the other distinct quadruples). In the notation of
\cref{lem:l20-star-mean} the three- and four-element supports therefore
contribute $\E_{\mu_p}[\mathcal W_3/p]=0$ and
$\E_{\mu_p}[\mathcal W_4/p]=(M_{G_5}+8q^2-33q^3+25q^4)/Z_4$
\eqref{eq:l20-star-identities}, and $|M_{G_5}|\le(1-q)(q^2+2qr)$
(\cref{lem:record:weighted-moment}(b)) gives $|m_{\mathsf K}-2q|\le m_\star/4$.

\emph{The Grams.} Each operator below commutes with the square-affine maps
acting on the label input and on all coordinates, so on $\mathcal H$ it is
$\alpha I+\gamma\Phat$, where $\alpha$ is its trace divided by $p-1$ and
$\gamma(p-1)$ is the trace of its product with $\Phat^*$; if it is
self-adjoint, $\gamma=0$ (\cref{lem:spec:affpair}(d)). For an orthonormal basis
$(h)$ of $\mathcal H$, $\sum_hh(y)h(y')=p\one_{y=y'}-1$, which gives the traces.

(a) Two core coordinates on the same side form a uniform distinct pair, and
$(a_i,c_j)$ a uniform directed pair, so $V_k^*V_{k'}=-\frac q{1-q}I$ for a
same-side pair and $V_{a_i}^*V_{c_j}=-\frac q{1-q}I+\frac r{1-q}\Phat$. The
coefficients of $\mathbf u$ sum to zero, and the Paley parts of
$V_{a_i}^*V_{c_j}$ and $V_{c_j}^*V_{a_i}$ cancel in $\mathbf u^*\mathbf u$ and
$\mathbf w^*\mathbf w$ and add up in $\mathbf u^*\mathbf w$; this gives (a) for
$\mathbf u$ and $\mathbf w$. The pair $(b_0,b_1)$ is a uniform distinct pair,
which gives $\mathbf v^*\mathbf v$.

(b) By \cref{lem:l20-star-matrix}(c),
$\E_z\mathsf K_\star^2=16\bar{\mathsf K}+1+4-2\bar{\mathsf K}-8\bar{\mathsf K}
=6\bar{\mathsf K}+5$ and $\E_z\mathsf K_\star=\bar{\mathsf K}-1$, so the trace of
$W_\star^*W_\star$ is $\E_{\mu_p}\Var_z\mathsf K_\star=11+6m_{\mathsf K}-v_{\mathsf K}$.
For a core coordinate $k$, $\mathsf K(x_k,C)=0$ and
$\ell(x_k,C)=-2\pm\chi(x_k-x_{k'})$, $k'$ the other coordinate on the same side,
so $\E_{\mu_p}\mathsf K_\star(x_k,C)=1$ by the exchange of $x_k$ and $x_{k'}$, and
the trace of $V_k^*W_\star$ is $1-m_{\mathsf K}$. Where $\mathsf K\ne0$,
$\chi(z-x_k)=s_k$, and $\E_z\chi(z-x_k)\ell=s_k$; so
$\E_z\chi(z-x_k)\mathsf K_\star=s_k(\bar{\mathsf K}-1)$ and
$V_k^*W_\star=\frac{q(1-m_{\mathsf K})}{1-q}I+\frac{s_km_{\mathsf K}r}{1-q}\Phat$.
Summing with the coefficients of $\mathbf u$ and of $\mathbf w$ gives (b).

(c) By (a), (b), $\Phat^2=-I$ and $\Phat^*=-\Phat$ on $\mathcal H$,
$15\,\mathbf z^*\mathbf z=W_\star^*W_\star/q+8m_{\mathsf K}(1-q)^{-1}I
+4\,\mathbf u^*\mathbf u$ and
$\sqrt{15}\,\mathbf w^*\mathbf z=\mathbf w^*W_\star/r-2\mathbf w^*\mathbf u\Phat$;
this gives $\mathbf z^*\mathbf z$ and $\mathbf w^*\mathbf z$, and
$\mathbf z^*\mathbf z\preceq g_{\mathbf z}I$ as $v_{\mathsf K}\ge0$. Given the
core, each pair coordinate is uniform on the $p-4$ labels outside it, so
$\E_{\mu_p}[h(b_l)\mid C]=-2(\mathbf wh)(C)/(p-4)$ for $h\in\mathcal H$; hence
$\mathbf v^*Y=-\frac{2\sqrt2\,q}{1-4q}\mathbf w^*Y$ for every operator $Y$ with
values in the functions of the core, and $Y=\mathbf w,\mathbf z$ give
$\mathbf v^*\mathbf w$ and $\mathbf v^*\mathbf z$.

(d) \Cref{lem:general-star}(a) with $w_{23}=-q$ gives $\mathbf e^*\mathbf e$.
The trace of $V_2^*W_{23}$ is $\E[\chi(0)\chi(b_0-b_1)]+q=q$, and the trace of
its product with $\Phat^*$ is $r^{-1}\E\,\E_z\chi(z-b_0)^2\chi(z-b_1)
=-qr^{-1}\E\chi(b_0-b_1)=0$; so
$V_2^*W_{23}=V_3^*W_{23}=\frac{q^2}{1-q}I$, which gives $\mathbf v^*\mathbf e$.
Given the core, $(b_0,b_1)$ is uniform on the ordered distinct pairs outside it,
and expanding $\sum_{y\ne y'}\chi(z-y)\chi(z-y')$ over such pairs gives, on
$\mathcal H$,
\[
 \E_{\mu_p}[W_{23}\mid C]=\frac{-2q\sum_kV_k+2\sum_{|S|=2}W_S}{(p-4)(p-5)},
\]
the sums over the core coordinates and over the six two-element subsets of the
core. As $\norm{V_k}=1$ and $\norm{W_S}\le r$ (\cref{lem:general-star}(a)), the
conditional expectation $\Pi_C\mathbf e$ of $\mathbf e$ onto the functions of the
core has norm at most $\rho$. For $Y=\mathbf w,\mathbf z$, which take values in
the functions of the core, $\mathbf e^*Y=(\Pi_C\mathbf e)^*Y$, and
$\norm{\mathbf w}\le1$, $\norm{\mathbf z}\le\sqrt{g_{\mathbf z}}$.

\emph{The consequence.} As $-\Delta_\star\preceq\Xi\preceq\Delta_\star$,
$\norm{\mathfrak F\Xi\mathfrak F^*}\le\norm{\mathfrak F\Delta_\star\mathfrak F^*}
=\norm{\Delta_\star^{1/2}\mathfrak G\Delta_\star^{1/2}}$, where
$\mathfrak G=\mathfrak F^*\mathfrak F$ is the Gram of the four features. In the
blocks $(\mathbf v,\mathbf z)$ and $(\mathbf w,\mathbf e)$ this matrix has
diagonal blocks of norms at most $k_{\rm top}$ and $k_{\rm low}$, by (a)--(d),
and an off-diagonal block of squared norm at most $k_\times^2$, the weighted sum
of the squared norms of its four entries. Since
$k_{\rm top}\ge\sqrt{30}>k_{\rm low}$, the largest eigenvalue of
$\bigl(\begin{smallmatrix}k_{\rm top}&k_\times\\k_\times&k_{\rm low}\end{smallmatrix}\bigr)$,
at most $k_{\rm top}+k_\times^2/(k_{\rm top}-k_{\rm low})\le k_4$, bounds the
norm. By \eqref{eq:l20-rank4-operator}, $T_\star/r=\mathfrak F\Xi_\star\mathfrak F^*$
with $\Xi_\star=\bigl(\begin{smallmatrix}0&\sqrt{30}\\\sqrt{30}&0\end{smallmatrix}\bigr)
\oplus\bigl(\begin{smallmatrix}0&2\\2&0\end{smallmatrix}\bigr)$, and
$-\Delta_\star\preceq\Xi_\star\preceq\Delta_\star$.
\end{proof}

The coefficient $k_4$ is $\sqrt{30}+O(q)$: the only overlap of order $r$, the
entry $\rho_{\mathbf z\mathbf w}$ between the two blocks, enters quadratically
across the gap $\sqrt{30}-2$.

\begin{proposition}[The soft L20 kernel as a swapped rectangle product]\label{prop:l20-rectangular}
Let $p\equiv3\pmod4$, $p\ge2^{20}$. Let $\mu^{\rm core}$ be the uniform law on the
admissible cores $C$ (distinct labels with $a_i\to c_j$) and $\mu^{\rm pair}$ the
uniform law on the ordered distinct pairs $b$, and let
$\mathsf R\colon L^2(\mu^{\rm pair})\to L^2(\mu^{\rm core})$ have the kernel
$\mathsf R(C;b)=\mathsf K(b_0,C)\mathsf K(b_1,C)$; the features $\mathbf z,\mathbf w$
of the core and $\mathbf v,\mathbf e$ of the pair are those of
\cref{lem:l20-star-matrix}, read in these laws.
\begin{enumerate}[label=\textup{(\alph*)}]
\item $\mathsf R$ has the row mass
 $\mathfrak r(C)=\mathsf R\one(C)=(\bar{\mathsf K}(C)^2-16q\bar{\mathsf K}(C))/(1-q)$
 and a constant column mass, and $\mathsf R_c=\mathsf R-\mathfrak r\otimes\one$
 (where $(\mathfrak r\otimes\one)h=(\E h)\,\mathfrak r$) satisfies
 $\mathsf R_c\one=\mathsf R_c^*\one=0$ and $\mathsf R_c^*\mathfrak r=0$.
\item Let $N_A$ be an upper bound for the norm of the double-centred
 nonstar remainder
 $P_0(\text{nonstars})P_0$. A raw cap from
 \cref{cor:soft-kernels}(a) suffices; the sharper bound of
 \cref{thm:rectangle-gram} also applies. Then
 \begin{equation}\label{eq:l20-rectangle}
  \mathsf R=\mathfrak r\otimes\one+\mathsf R_\star+\mathsf R_\circ,\qquad
  \mathsf R_\star=r\bigl(\sqrt{30}\,\mathbf z\mathbf v^*+2\mathbf w\mathbf e^*\bigr),\qquad
  \norm{\mathsf R_\circ}\le N_A ,
 \end{equation}
 and $\norm{\mathsf R_\star}\le k_4r$; hence $\norm{\mathsf R_c}\le k_4r+N_A$.
\item On $L^2(\mu^{\rm core}\otimes\mu^{\rm pair})$ the soft L20 kernel is
 \begin{equation}\label{eq:l20-swap}
  K_{\rm soft}\bigl((C,b),(C',b')\bigr)=\mathsf R(C;b')\,\mathsf R(C';b),\qquad
  K_{\rm soft}=T[\mathsf R,\mathsf R]=(\mathsf R\otimes\mathsf R^*)\operatorname{Swap}.
 \end{equation}
\item Put $\eta=2q+11r$, $h_p=(1-4q)(1-5q)/(1-q)$ and
 \[
  s_p=\frac{\bigl[1+4\mathfrak m_++(6+4\eta+\eta^2)\eta^2
  +256q^2(1+2\mathfrak m_++\eta^2)\bigr]^{1/2}}{1-q}.
 \]
 Then $\norm{\mathfrak r}\le s_p$. If $\norm{\mathsf R_c}\le c_{\rm R}<1$, then for
 every collision cap $D_{\rm col}$ with $\eps_Z=N_{\rm ns}+D_{\rm col}+m_\star<1$
 as in \cref{prop:l20-completion},
 \[
  T\succeq-C'I,\qquad\norm{T|_{\mathcal I^\perp}}\le C',\qquad
  C'=\frac{s_pc_{\rm R}/h_p+D_{\rm col}}{1-\eps_Z}.
 \]
\end{enumerate}
\end{proposition}

\begin{proof}
(a) Given $C$, under $\mu^{\rm pair}$, $\mathfrak r(C)=\frac1{p(p-1)}\sum_{y\ne y'}
\mathsf K(y,C)\mathsf K(y',C)=\frac{p^2\bar{\mathsf K}^2-\sum_y\mathsf K(y,C)^2}{p(p-1)}$,
which is the stated formula by $\mathsf K^2=16\mathsf K$. If $\varphi$ is a
function of the core invariant under the square-affine maps, then
$\mathsf R^*\varphi(b)=\E_{\mu^{\rm core}}[\mathsf R(C;b)\varphi(C)]$ is invariant
under these maps acting on $b$, hence constant on each of the two orbits of
ordered distinct pairs (the sign of $\chi(b_1-b_0)$); $\mathsf R$ is symmetric
in $b_0,b_1$, and the exchange maps one orbit onto the other because
$\chi(-1)=-1$, so $\mathsf R^*\varphi$ is constant. With $\varphi=\one$ and
$\varphi=\mathfrak r$ this gives the constant column mass $\E\mathfrak r$ and
$\mathsf R^*\mathfrak r=\norm{\mathfrak r}^2\one$; hence $\mathsf R_c\one=0$,
$\mathsf R_c^*\one=0$ and $\mathsf R_c^*\mathfrak r=0$.

(b) By (a), $\mathsf R_c=P_0\mathsf RP_0$, with $P_0$ the projections onto the
mean-zero functions. Expand
$\mathsf R=\prod(1+s_k\chi(b_l-x_k))$ over the eight cross pairs (core
coordinate $k$, pair coordinate $b_l$) into its $256$ subset terms: the
constant, $34$ stars ($8$ single pairs, $22$ sets of at least two pairs at one
pair coordinate and $4$ pairs $\{(k,b_0),(k,b_1)\}$ at one core coordinate)
and $221$ nonstars. The double centring removes the constant. The stars at
$b_l$ add up to $(\ell+\mathsf K_\star)(b_l,C)$, and
$\E_z\ell(z,C)h(z)=\sum_ks_k(Ph)(x_k)=-2r(\mathbf u\Phat h)(C)$; so after the
double centring they give $(W_\star-2r\mathbf u\Phat)V^*_{b_l}=\sqrt{15}\,r\,\mathbf zV^*_{b_l}$,
with $V_{b_l}h=h(b_l)$; the stars at the core coordinates give
$\sum_kV_kW_{23}^*=2r\mathbf w\mathbf e^*$, their rank-one marginal terms being
exactly those removed by the double centring (\cref{cor:star-matrix}).
Hence the centred stars give
$\sqrt{15}\,r\,\mathbf z(V^*_{b_0}+V^*_{b_1})+2r\mathbf w\mathbf e^*=\mathsf R_\star$,
and $\mathsf R_\circ=P_0(\text{nonstars})P_0$ has norm at most that of the sum
of the nonstar terms. The raw extraction bound of \cref{cor:soft-kernels}(a)
also bounds this compression: the core and pair coordinate embeddings
are isometries, and the eight-edge kernel uses the same $221$ subsets.
The sharper direct calculation is \cref{thm:rectangle-gram}.
Finally $\mathsf R_\star$ maps the mean-zero functions of the pair to those of
the core. Its Hermitian dilation on $L^2(\mu^{\rm core})\oplus L^2(\mu^{\rm pair})$ is
$r\mathfrak F'\Xi_\star\mathfrak F'^*$ for the four features placed in their
summands; their Gram has the blocks of \cref{prop:l20-rank4} inside each summand
($\mathbf z,\mathbf w$ depend only on the core and $\mathbf v,\mathbf e$ only on
the pair, whose laws under $\mu_p$ are $\mu^{\rm core}$ and $\mu^{\rm pair}$) and
zero blocks across them. The argument of \cref{prop:l20-rank4} gives
$\norm{\mathsf R_\star}\le k_4r$.

(c) The cross factors between the core of $x$ and the pair of $x'$ are
$1-\chi(b'_l-a_i)$ and $1+\chi(b'_l-c_j)$, the entries $-1$ and $1$ of $D_{20}$
in the rows of $a_i$ and $c_j$; their product is $\mathsf R(C;b')$. By the
antisymmetry of $D_{20}$ the factors between the pair of $x$ and the core of
$x'$ are $1+\chi(a'_i-b_l)=1-\chi(b_l-a'_i)$ and $1-\chi(c'_j-b_l)=1+\chi(b_l-c'_j)$,
with product $\mathsf R(C';b)$, and $D_{20}$ has no other nonzero entries. The
second form is \cref{lem:swapped-rectangle}(a).

(d) Write $\bar{\mathsf K}=1+y$. By \cref{lem:l20-star-matrix}(c) and
\cref{thm:weil}, $|y|\le2q+4\cdot2r+3r=\eta$ (six two-root, four three-root and
one four-root sum), and $\E_{\mu^{\rm core}}y=m_{\mathsf K}\le\mathfrak m_+$
(\cref{prop:l20-rank4}). Hence
$\E\bar{\mathsf K}^4\le1+4\mathfrak m_++(6+4\eta+\eta^2)\eta^2$ and
$\E\bar{\mathsf K}^2\le1+2\mathfrak m_++\eta^2$, and dropping the term
$-32q\bar{\mathsf K}^3\le0$ in $(1-q)^2\mathfrak r^2$ gives
$\norm{\mathfrak r}\le s_p$; note $s_p\ge1$. By (a), $\mathsf R$ has the singular
pair $(\mathfrak r/\norm{\mathfrak r},\one)$ with the value $\norm{\mathfrak r}$,
and all its other singular values are at most $c_{\rm R}$. By
\cref{lem:swapped-rectangle}(a) the eigenvalues of $T[\mathsf R,\mathsf R]$ are
$0$, the squares of the singular values and their products $\pm\sigma_i\sigma_j$
($i\ne j$); as $c_{\rm R}<1\le s_p$, all negative eigenvalues are at least
$-s_pc_{\rm R}$, and all eigenvalues except the one at the translation-invariant
vector $(\mathfrak r/\norm{\mathfrak r})\otimes\one$ have modulus at most
$s_pc_{\rm R}$. The law $\mu_p$ is $\mu^{\rm core}\otimes\mu^{\rm pair}$ conditioned on
the translation-invariant event $\mathcal E$ that the pair avoids the core,
which has probability $h_p$ for every core. With the isometry
$Jf=\one_{\mathcal E}f/\sqrt{h_p}$ of $L^2(\mu_p)$ into
$L^2(\mu^{\rm core}\otimes\mu^{\rm pair})$, $K_{\rm soft}=h_p^{-1}J^*T[\mathsf R,\mathsf R]J$
on $L^2(\mu_p)$, and $J$ maps $\mathcal I^\perp$ into the orthogonal complement
of the translation-invariant functions. So $K_{\rm soft}\succeq-(s_pc_{\rm R}/h_p)I$
and $\norm{K_{\rm soft}|_{\mathcal I^\perp}}\le s_pc_{\rm R}/h_p$. Finally
$T=(K_{\rm soft}-T_{\rm coll})/Z_0$ with $\norm{T_{\rm coll}}\le D_{\rm col}$ and
$Z_0\ge1-\eps_Z$ (proof of \cref{prop:l20-completion}).
\end{proof}

\subsubsection{Bounding the centred remainder through its Gram}

The $221$ nonstar terms of the rectangle share their four core
variables. Estimating the Gram of their sum retains cancellations
that are lost when the $221$ terms are bounded separately and their
norms are then added. The following estimate
is for precisely $\mathsf R_\circ$ in \eqref{eq:l20-rectangle}.

\begin{lemma}[Two-sided bounds for the squared cubic trace $H_u$]\label{lem:rectangle-mellin}
For $u\in\Fp$, on counting measure on $\Fp\setminus\{u\}$, put
\[
 T_u(v,w)=\E_y\chi(u-y)\chi(v-y)\chi(w-y),\qquad H_u=T_u\circ T_u.
\]
Then $\|T_u\|\le1$, $\|H_u\|\le1-2q-2q^2$, and
\[
 H_u\succeq-(3r+q-q^2)I.
\]
\end{lemma}
\begin{proof}
The full character matrix $C$ satisfies $CC^*=pI-J$.
The matrix $T_u$ is a compression of
$qC\diag(\chi(u-y))C^*$, so has norm at most $1$.
Character orthogonality gives each row sum of the nonnegative matrix
$H_u$ as $1-2q-2q^2$, proving its absolute norm bound.
For the lower bound translate $u$ to zero. Dilation gives
$T_0(v,w)=\chi(v)T_0(1,w/v)$, so $H_0$ is multiplicative convolution.
Its trivial-character eigenvalue is its positive row sum. For a
nontrivial multiplicative character $\rho$, the eigenvalue is
\[
 -q+q^2+p^{-2}\sum_{\lambda\ne0,1}\rho(\lambda)
  \tr(\operatorname{Frob}_\lambda\mid\operatorname{Sym}^2\mathcal F),
\]
where $\mathcal F$ is the Legendre sheaf of
\cref{thm:prelim:ks} with $h(X)=X(X-1)$. Indeed its determinant is
$\overline{\Q}_\ell(-1)$, so the symmetric-square trace is the square
of the cubic character sum minus $p$; the missing value at $1$ is $q^2$.
The Kummer sheaf for $\rho$ is the corresponding eigensheaf of
$z\mapsto z^{p-1}$ on $\mathbb G_m$. It has weight zero and is tame.
The twist of $\operatorname{Sym}^2\mathcal F$ has rank $3$, weight $2$,
and is tame at $0,1,\infty$. Its geometric irreducibility follows from
the $\mathrm{SL}_2$ monodromy of $\mathcal F$ and irreducibility of
the symmetric square. Thus its invariants and coinvariants vanish.
By \cref{thm:prelim:curves,thm:prelim:gos}, $H_c^0=H_c^2=0$ and
$\dim H_c^1=3$. The trace formula and Deligne's bound
(\cref{thm:prelim:lefschetz,thm:prelim:weights}) bound the sum by
$3p^{3/2}$. All eigenvalues are real because $H_0$ is symmetric.
This proves the stated lower bound.
\end{proof}

\begin{theorem}[The centred accepting rectangle]\label{thm:rectangle-gram}
For $p\equiv3\pmod4$, $p\ge2^{20}$, let
$Z_4=(1-q)(1-3q)(1-7q)$. Then
\begin{equation}\label{eq:rectangle-gram-cap}
 \|\mathsf R_\circ\|\le
 4q\sqrt{\frac{920+4012r+48641q+3008q^2+320q^3}{(1-q)Z_4}}
 \le121.588q.
\end{equation}
For $p\ge2^{36}$ the last bound improves to $121.328q$.
Without \cref{lem:rectangle-mellin}, the same argument gives $128q$.
These are centred bounds; they do not replace a raw nonstar or
collision bound elsewhere in the proof.
\end{theorem}

\begin{proof}
Put $h=1-q$ for this proof. Initially give the four core coordinates
independent uniform laws and retain the uniform distinct private pair.
With $s=(-1,-1,1,1)$, set
\[
 k_x(t)=\prod_{i=1}^4(1+s_i\chi(t-x_i)),\qquad
 N(x;b_0,b_1)=(k_x(b_0)-1)(k_x(b_1)-1)
       -\sum_i\chi(b_0-x_i)\chi(b_1-x_i).
\]
This is exactly the sum of the nonstar terms. The accepting-core
event has probability $Z_4/16$: each distinct ordered $c$-pair has
$(p-3)/4$ common predecessors, from which the ordered $a$-pair is
chosen distinctly. Restriction therefore costs $4/\sqrt{Z_4}$;
output centring is a contraction. We bound $NP_0$ on the full core cube.

\emph{The Gram away from private overlaps.}
For four distinct private labels $z_1,z_2,z_3,z_4$, write
\[
 v=\E_y\prod_i\chi(z_i-y),\qquad
 t_i=\E_y\prod_{j\ne i}\chi(z_j-y),\qquad T=\sum_i t_i.
\]
Independent integration in the four core variables gives
\begin{align}
 G={}&[(1-6q+v)^2-T^2]^2
       -\sum_i[(1-3q)^2-t_i^2]^2+6h^4-3\notag\\
    &-8[h^3(v-q)+q]+4v+12q^2.\label{eq:rectangle-gram-polynomial}
\end{align}
Denote the same polynomial on all private pair pairs by $\widetilde G$.
With $a=1-6q$, $b=1-3q$, its nonconstant part is
\begin{align*}
 \alpha v+6a^2v^2+(12q-54q^2)\sum_i t_i^2
 -4a^2\sum_{i<j}t_it_j+4a(v^3-vT^2)+Q_4,\\
 \alpha=-48q+408q^2-856q^3<0,\qquad
 Q_4=v^4-2v^2T^2+T^4-\sum_i t_i^4.
\end{align*}
Powers denote pointwise products. The full kernel $v$ is a positive
Gram, so the term with negative coefficient is $\preceq0$ in Loewner order.
The constant kernel vanishes after compression by $P_0$.

\emph{Quadratic terms.}
Put $g=(1-2q)/(1-q)$. Applying the four-cycle estimate to the
off-diagonal output of the pair-to-field-square map, and the singleton
incidence map to its diagonal, gives
\[
 \|P_0v^2P_0\|\le q^2\left(\frac4h+2gq\right).
\]
For $t_it_j$ with both omissions on the same private half, two outer
Paley transforms give
$q^2\sqrt{4/h^2+2gq/h}$; there are two such pairs.
For each of the four opposite-half pairs, the middle kernel is a
two-edge matching, giving $q^2/h$. These bounds follow from
\cref{cor:dl:normalized,cor:raw-pair-maps,lem:matching:kernel}
by keeping the common auxiliary labels. Together they bound the
quadratic cross terms by $(56/h+16q)q^2$.
For $t_i^2$, split its two auxiliary labels into distinct and equal
ones. The distinct part factors through the raw pair-to-singleton
map and the four-cycle map and has norm at most $2q^{3/2}$.
The equal part is
$q(1-3q)J+q^2(\one[c=d]+\one[c=e])$ on the singleton/pair chart.
After centring its norm is at most $\sqrt{2g}q^3$. Thus the four
square terms cost at most $(96r+96q^2)q^2$.

\emph{Cubic auxiliary collisions.}
For four distinct private labels put
$A(y)=\prod_i\chi(z_i-y)$ and
$B(y)=\sum_i\prod_{j\ne i}\chi(z_j-y)$.
Let $D_3$ be the unrenormalized average of
$A(u)A(v')A(w)-A(u)B(v')B(w)$ over three distinct auxiliary labels.
The equality-partition formula for three summation labels gives
\[
 v^3-vT^2=D_3-qv+q^2(6+L)v-2q^3L,
 \qquad L=\sum_i\left(\sum_{j\ne i}\chi(z_j-z_i)\right)^2\le20.
\]
For verification, the merged moments are
$\E A^2=1-4q$, $\E A^3=v$, $\E AB=0$,
$\E B^2=4(1-3q)-qL$, and $\E AB^2=4v-qL$.
The bound on $L$ is the score bound for a four-vertex tournament.
Since $|v|\le2r+q$, the error away from overlaps is at most
$52q^2r+66q^3$.
On the private diagonal the error after $-qv$ is
$-4q+40q^2-80q^3$, on a row fraction $2q^2/h$.
On exactly one overlap it is at most $73q^2$, on a row fraction
at most $4q/h$. These follow by the same three-label partition:
there $|v|\le2q$, $|T|\le4r+2q$, $|\E AB|\le2q$,
$\E A^2\le1$, $|\E B^2|,|\E AB^2|\le16$.
The absolute row and column bounds therefore give
\[
 v^3-vT^2=D_3-qv+E_3,\qquad
 \|E_3\|\le52q^2r+400q^3/h.
\]
The term $D_3$ has two-label witness mass $28$: $4$ for $v^3$,
$8\cdot2$ for the ordered same-half omissions, and $8$ for the
opposite-half omissions. Matching, path and four-cycle bounds give
$\|D_3\|\le28q^2/h$. Auxiliary distinctness is an output
projection and is not renormalized; each private-half restriction
costs $h^{-1/2}$. Dropping the negative semidefinite term $-qv$, the cubic
contribution is at most $(112/h+208r+1600q/h)q^2$.

\emph{Quartic auxiliary collisions.}
For four support characters $\psi_i$, write $\tau_i=\E\psi_i$ and
$m_I=\E\prod_{i\in I}\psi_i$. The exact partition identity is
\begin{align*}
 \tau_1\tau_2\tau_3\tau_4={}&D_4
 +q\sum_{ij}m_{ij}\tau_k\tau_l
 -q^2\sum_{ij|kl}m_{ij}m_{kl}\\
 &-2q^2\sum_{ijk}m_{ijk}\tau_l+6q^3m_{1234},
\end{align*}
where the multiplicities are $6,3,4,1$ and $D_4$ retains four
distinct auxiliary labels. This identity includes every root-hit hole.
For any of the five traces $v,t_1,\ldots,t_4$, and any nonempty
XOR of at most three of their supports, the row and column second
moments are at most $q$, by Paley orthogonality and the raw pair
Gram. Hence a product of two such traces has absolute Schur norm
at most $q$. Also
\[
 \|P_0(\tau\tau')P_0\|\le2q^{3/2}/h+16q^2/h.
\]
The square-cubic case was proved above; different cubics use the
quadratic cross bounds. For $v$ times a cubic the distinct term
has a path and a rectangle witness, and the equal term has a
singleton trace with at most three holes, giving $5q^2/h$.
Each $m_I$ differs from its XOR trace by at most $4q$.
These observations give the following bounds after centring:
\[
\begin{array}{c|c}
 \text{one partition term}&\text{upper bound}\\ \hline
 qm_{ij}\tau_k\tau_l&2q^{5/2}/h+20q^3/h\\
 q^2m_{ij}m_{kl}&9q^3\\
 2q^2m_{ijk}\tau_l&3q^3\\
 6q^3m_{1234}&6q^3
\end{array}
\]
For the first row, identical supports use the preceding centred
bound, with $4q^3$ for holes. Nonidentical supports have singleton
or two-root XOR. The latter is $-q$ plus a possible cross-half
equality; on that equality the product of traces is at most $4q$,
apart from a diagonal of row mass $q^2/h$. This gives at most
$10q^3/h$ including holes. For the second row two identical
blocks give $q^2((1-qu_i)(1-qu_j)-1)$, $0\le u_i,u_j\le4$;
one or no identical blocks give respectively $6q^3$ or $2q^3$.
The third row has nonempty XOR and costs $2(1+4r)q^3$ by
row and column Cauchy--Schwarz. The last row uses $|m_I|\le1$.

The absolute word mass in $Q_4$ is $1+32+252=285$.
Its distinct-label witness mass is $332$: $4$ from $v^4$,
$2(8\cdot2+8)=48$ from $-2v^2T^2$, and
$28\cdot2+224=280$ from $T^4-\sum t_i^4$.
The same two-label witnesses as above therefore give
\[
 \|P_0Q_4P_0\|\le(332+3420r+47025q)q^2/h.
\]
Here $3420=285\cdot6\cdot2$ and
$47025=285(6\cdot20+3\cdot9+4\cdot3+6)$.

\emph{Private overlaps.}
It remains to bound $\Delta=N^*N-\widetilde G$. Use unordered
private pairs, of cardinality $n=\binom p2$; symmetry makes this
equivalent to the ordered-pair operator. A universal expression,
valid also at repeated roots, is obtained by setting
\[
 m_{\epsilon,J}=\E_y\prod_{j\in J}(1+\epsilon\chi(z_j-y)),\qquad
 F(z)=\sum_{J\subseteq[4]}(-1)^{4-|J|}m_{-,J}^2m_{+,J}^2.
\]
The true Gram is $F(z)-8[h^3(v-q)+q]+4v+12q^2$.
Subtracting \eqref{eq:rectangle-gram-polynomial} gives
\[
 \Delta_{A,A}=210-1512q+4020q^2-3720q^3.
\]
For $A=\{u,v_0\}$, $B=\{u,w\}$ put
$e=\chi(u-v_0)$, $f=\chi(u-w)$, $s=e+f$, $\eta=ef$, and
$t=T_u(v_0,w)$. The other nonzero entries are
\begin{align*}
 \Delta_{A,B}={}&-28(2+\eta)q-24t^2+(594+240\eta)q^2
 -(1204+528\eta)q^3\\
 &+(112+16\eta)qt^2+s(-24q+128q^2)t.
\end{align*}
These formulas follow by expanding the displayed $m_{\epsilon,J}$;
no distinct-root formula has been applied on an overlap.

For clarity, the one-sided operator bound can be checked with
incidence matrices. On counting measure let
$U(A,u)=\one[u\in A]$ and
$W(\{u,v_0\},u)=\chi(u-v_0)$, zero otherwise.
Their Grams have diagonal $2$ and off-diagonal entries $1$ and
$\eta$ on a shared endpoint. Their coefficients in $\Delta$ are
$-56q+594q^2-1204q^3$ and $-28q+240q^2-528q^3$, both negative.
After their diagonal terms are removed, their sum is at most
$(168q-1668q^2+3464q^3)I$.
Let $R_u$ restrict to pairs containing $u$, and let
$E_u=\diag_{v_0\ne u}\chi(u-v_0)$. Since $\sum_uR_u^*R_u=2I$,
the remaining trace-square and trace-linear terms are
\begin{align*}
 &(-24+112q)(H^{(0)}-2q^2I)+16q(H^{(1)}-2q^2I)\notag\\
 &\hspace{25mm}+(-24q+128q^2)(K+4qI),\\
 H^{(0)}&=\sum_uR_u^*H_uR_u,\quad
 H^{(1)}=\sum_uR_u^*E_uH_uE_uR_u,\quad
 K=\sum_uR_u^*(E_uT_u+T_uE_u)R_u.
\end{align*}
By \cref{lem:rectangle-mellin},
$H^{(0)}\succeq-2(3r+q-q^2)I$,
$\|H^{(1)}\|\le2(1-2q-2q^2)$, and $\|K\|\le4$.
Adding the diagonal and incidence contributions, and dividing by
$n$, gives the probability-normalized bound
\[
 \Delta\preceq\frac{q^2}{h}
 [420+(288-1344q)r-2336q+2912q^2+320q^3]I.
\]
Using only $\|H_u\|\le1-2q-2q^2$ instead gives
$\Delta\preceq516q^2I$ for $q\le1/16$.

\emph{Assembly.}
Add the quadratic, cubic, quartic and overlap estimates, discarding
only negative terms. On mean-zero private inputs this yields
\[
 \|NP_0\|^2\le
 q^2\frac{920+4012r+48641q+3008q^2+320q^3}{h}.
\]
The accepting-core restriction gives \eqref{eq:rectangle-gram-cap}.
The factor multiplying $q$ in the first bound is increasing in
$q$ and $r$ while $Z_4>0$.
Exact evaluation at $(q,r)=(2^{-20},2^{-10})$ and
$(2^{-36},2^{-18})$ gives respectively coefficients less than
$121.588$ and $121.328$. The elementary overlap bound gives instead
\[
 \|NP_0\|^2\le q^2\left(516+
            \frac{500+3724r+48641q+96q^2}{h}\right),
\]
which proves the stated $128q$ alternative on the same domain.
\end{proof}

The coefficient in \eqref{eq:rectangle-gram-cap} may be used as
$N_A/q$ in the gap estimate below. The selected completion weights
and their certificates retain the larger, previously certified
rectangle prices. They remain valid upper assignments; no global
saving is claimed merely by reducing this local norm bound.


\begin{proposition}[Sharper bound for the centred rectangle, using the gap $30q-\ell_g^2$]\label{prop:l20-pair-gap}
In the setting of \cref{prop:l20-rectangular} put
\begin{gather*}
 g_{\mathbf v}=\frac{1-2q}{1-q},\qquad
 M_g=15+44\mathfrak m_++(42+12\eta)\eta^2,\qquad
 \tau_g^2=\frac{2qM_g}{(1-q)^2(1-2q)},\\
 e_g=\frac{\sqrt{30}\,r(\eta+15q)\sqrt{g_{\mathbf z}}}{(1-q)\sqrt{g_{\mathbf v}}},\qquad
 w_g=\frac{\sqrt{30}\,r\rho_{\mathbf z\mathbf w}}{\sqrt{g_{\mathbf v}}}+e_g,\qquad
 l_g=2r+N_A,\\
 x_g=2rw_g+\tau_gN_A,\qquad\Delta_g=30q-l_g^2 .
\end{gather*}
If $\Delta_g>0$, then
$\norm{\mathsf R_c}\le c_{\rm gap}:=\bigl(\tau_g^2+x_g^2/\Delta_g\bigr)^{1/2}$.
\end{proposition}

\begin{proof}
On $\mathcal H$, $\mathbf v^*\mathbf v=g_{\mathbf v}I$, so
$\mathbf v'=\mathbf v/\sqrt{g_{\mathbf v}}$ is an isometry into the mean-zero
functions of the pair. As $\E_z\ell(z,C)h(z)=-2r(\mathbf u\Phat h)(C)$ (proof of
\cref{prop:l20-rectangular}(b)), $\mathbf zh=\E_z\mathsf K(z,C)h(z)/(\sqrt{15}\,r)$
for $h\in\mathcal H$. Given $b_0=y$, the label $b_1$ is uniform on the other $p-1$
labels and $\sum_{y'\ne y}\mathsf K(y',C)=p\bar{\mathsf K}-\mathsf K(y,C)$; with
$\mathsf K^2=16\mathsf K$ this gives the exact action
\[
 \mathsf R\mathbf v=\sqrt{30}\,r\,\frac{\bar{\mathsf K}-16q}{1-q}\,\mathbf z .
\]
The Gram $(\mathsf R\mathbf v')^*\mathsf R\mathbf v'$ is self-adjoint and commutes
with the square-affine maps, so it is a scalar (\cref{lem:spec:affpair}(d)); its
trace is $30q\,\E[(\bar{\mathsf K}-16q)^2(1-q)^{-2}\norm{P_0\mathsf K(\cdot,C)}^2]
/(15q\,g_{\mathbf v})$ with $\norm{P_0\mathsf K(\cdot,C)}^2=16\bar{\mathsf K}-\bar{\mathsf K}^2$,
and division by $p-1$ gives
\[
 (\mathsf R\mathbf v')^*\mathsf R\mathbf v'
 =\frac{2q\,\E\bigl[(\bar{\mathsf K}-16q)^2(16\bar{\mathsf K}-\bar{\mathsf K}^2)\bigr]}
 {(1-q)^2(1-2q)}I .
\]
Here $\bar{\mathsf K}=16n/p=16nq$, where $n=\#\{z:\mathsf K(z,C)=16\}$; the
weight vanishes at $n=0$,
and for $n\ge1$, $0\le\bar{\mathsf K}-16q\le\bar{\mathsf K}\le16$, so the
expectation is at most $\E(16\bar{\mathsf K}^3-\bar{\mathsf K}^4)$; expanding at
$\bar{\mathsf K}=1+y$, $16\bar{\mathsf K}^3-\bar{\mathsf K}^4=15+44y+42y^2+12y^3-y^4$,
and $|y|\le\eta$, $\E y\le\mathfrak m_+$ bound it by $M_g$. Thus
$\norm{\mathsf R\mathbf v'}\le\tau_g$, and $\tau_g^2\ge30q$ as $M_g\ge15$. On the
orthogonal complement of the range of $\mathbf v'$ in the mean-zero functions of
the pair, $\mathbf v^*=0$, so there $\mathsf R_c=2r\mathbf w\mathbf e^*+\mathsf R_\circ$
has norm at most $l_g$, as $\norm{\mathbf w},\norm{\mathbf e}\le1$. Moreover
$|(\bar{\mathsf K}-16q)/(1-q)-1|\le(\eta+15q)/(1-q)$ gives
$\norm{\mathsf R\mathbf v'-\sqrt{30}\,r\mathbf z/\sqrt{g_{\mathbf v}}}\le e_g$ and
$\norm{(\mathsf R\mathbf v')^*\mathbf w}\le w_g$, so the off-diagonal block of
$\mathsf R_c^*\mathsf R_c$ in this decomposition has norm at most
$2rw_g+\tau_gN_A=x_g$. The diagonal blocks have norms at most $\tau_g^2$ and
$l_g^2$, and $\tau_g^2-l_g^2\ge\Delta_g>0$, so the largest eigenvalue of the
scalar comparison matrix, and with it $\norm{\mathsf R_c}^2$, is at most
$\tau_g^2+x_g^2/\Delta_g$.
\end{proof}

The nonstar term $N_A$ enters $c_{\rm gap}^2$ through $x_g$, of order $qr$,
across the gap $\Delta_g$ of order $q$, and through $l_g$ in the gap itself, at
relative order $r$; so it changes $(c_{\rm gap}/r)^2$ by $O(q)$, whereas in
$k_4r+N_A$ it enters $c/r$ linearly, at order $r$.

\begin{proposition}[Collisions with the core--pair weights of the proof; the caps of L20]
\label{prop:l20-collisions}
For $p\equiv3\pmod4$, $p\ge2^{20}$, put $E_{c,j}=e_j(p,12)
=((j/2-1)2^j+1)r+12\cdot2^jq$ ($j=2,3,4$),
$\Lambda_c=(1-e_2(p,6))^4(1-e_0(p,6))^2=(1-r-24q)^4(1-6q)^2$,
$U_c=(1+E_{c,4})^6$,
\begin{gather*}
 P_c(q)=96q+3712q^2+81920q^3+1245184q^4+12582912q^5+67108864q^6,\\
 S_c=32q(1+E_{c,3})^2(1+E_{c,4})^3+64q(1+E_{c,2})(1+E_{c,3})^2(1+E_{c,4})^2,\\
 D_{\rm bal}=\frac{S_c+U_c(P_c-96q)}{\Lambda_c}.
\end{gather*}
\begin{enumerate}[label=\textup{(\alph*)}]
\item $\norm{T_{\rm coll}}\le D_{\rm bal}$ for the soft L20 kernel.
\item Put $\eps_Z=N_{\rm ns}+D_{\rm bal}+m_\star$,
 $C_{20}=(s_pc_{\rm gap}/h_p+D_{\rm bal})/(1-\eps_Z)$,
 $H_{20}=(N_{\rm ns}+D_{\rm bal})/(1-\eps_Z)$, $g_{20}=C_{20}/r$ and
 $c^+_{20}=p\chi^+_p/(1-H_{20})$. If $\eps_Z<1/2$, $\Delta_g>0$,
 $c_{\rm gap}<1$, $C_{20}<1$ and $H_{20}<1$ at an endpoint $P\ge2^{20}$ (that
 is, with $q=1/P$ and $r=P^{-1/2}$), then
 $-g_{20}(P)p^{-1/2}\le\mathsf A_p\le1+c^+_{20}(P)/p$ for every prime $p\ge P$
 with $p\equiv3\pmod4$, and a middle block $\Blk$ of a piece starting at
 $X\ge P$, with least prime $b$, has
 \[
  \mathsf A_{\Blk}\succeq-g_{20}(P)\exp\bigl(c^+_{20}(P)L_{\rm loc}/(X\log X)\bigr)
  \,b^{-1/2}.
 \]
\end{enumerate}
\end{proposition}

\begin{proof}
(a) The soft kernel contains the factors $\mathsf K(b'_l,C)$ and
$\mathsf K(b_l,C')$ (\cref{prop:l20-rectangular}(c)), which vanish at the core
labels (\cref{lem:l20-star-matrix}(c)). So the strata of $T_{\rm coll}$ with a
coincidence between a pair label of one half and a core label of the other
contribute nothing, and in the bounds of \cref{cor:incident-weights} the weights
$\alpha_{ij}(x)$ of the $16$ pairs $(i,j)$ joining a core coordinate and a pair
coordinate may be replaced by zero: a stratum containing such a pair is bounded
by zero, and the bounds of the other strata do not change. The remaining
coincidences $x'_i=x_j$ join two core coordinates ($16$ choices of $(i,j)$) or two
pair coordinates ($4$ choices); for these, $j\notin N(i)$ and $\mathrm{mg}(j,i)=0$
in \cref{lem:oct24-collision}. Fix $x$ and apply \cref{cor:incident-weights} with
the actual incident weights $\alpha_{ij}(x)=w_i(x_j)$, where $w_i$ is the
cross factor at vertex $i$ of that corollary. For an $a$-coordinate
$i$ of $x'$ and a core coordinate $j$,
$w_i(x_j)=\prod_{l=0,1}(1-\chi(b_l-x_j))$; for a $c$-coordinate,
$w_i(x_j)=\prod_{l=0,1}(1+\chi(b_l-x_j))$; for a pair coordinate $i$ and a pair
coordinate $j$, $w_i(x_j)=\mathsf K(x_j,C)$. At a core label $y$ of $x$ the
labels $y,b_0,b_1$ are distinct, so the two products are $0$ or $4$, and not
both are $4$. Let $n_-$ and $n_+$ count the core labels of $x$ at which the
first, respectively the second, product is $4$; then $n_-+n_+\le4$, and the row
charges $\rho_i(x)=q\sum_j\alpha_{ij}(x)$ are $\rho_a=4n_-q$ for the two
$a$-coordinates, $\rho_c=4n_+q$ for the two $c$-coordinates and
$\rho_{b_l}=q(\mathsf K(b_0,C)+\mathsf K(b_1,C))\le32q$ for the two pair
coordinates.

\emph{One coincidence.} Order the other vertices so that, for a coincident core
coordinate, its two neighbours on the opposite side of the core come before the
other coordinate on its side, and, for a coincident pair coordinate, the core
alternates between the sides. The five remaining vertices then have
$t^{ij}_\ell$ equal to $3,3,4,4,4$, respectively $2,3,3,4,4$, and
\cref{cor:incident-weights} bounds the row sum of $T_{\rm coll,1}$ by
$\Lambda_c^{-1}$ times
\[
 \sum_{i\ \rm core}\rho_i(x)(1+E_{c,3})^2(1+E_{c,4})^3
 +\sum_{i\ \rm pair}\rho_i(x)(1+E_{c,2})(1+E_{c,3})^2(1+E_{c,4})^2\le S_c,
\]
since the core charges add up to $8(n_-+n_+)q\le32q$ and the pair charges to at
most $64q$; here $\Lambda_c=L_{\rm int}$ and $e_t(p,2n)=E_{c,t}$.

\emph{Two or more coincidences.} Every vertex has degree $d_i=4$, so
$\Upsilon_k\le(1+E_{c,4})^6=U_c$, and the row sum of $T_{\rm coll,\ge2}$ is at
most $U_c\Lambda_c^{-1}$ times the part of degree at least two of
$\sum_I2^{e(I)}\prod_{i\in I}\rho_i(x)$. The internal edges join the
$a$-coordinates to the $c$-coordinates, so this sum is
$F(\rho_a,\rho_c)(1+\rho_{b_0})(1+\rho_{b_1})$ with
\begin{align*}
 F(\rho_a,\rho_c)&=\sum_{k,l=0}^2\tbinom2k\tbinom2l2^{kl}\rho_a^k\rho_c^l\\
 &=1+2(\rho_a+\rho_c)+(\rho_a+\rho_c)^2+6\rho_a\rho_c
 +8(\rho_a+\rho_c)\rho_a\rho_c+16\rho_a^2\rho_c^2 .
\end{align*}
As a polynomial in $\rho_a+\rho_c$, $\rho_a\rho_c$, $\rho_{b_0}$, $\rho_{b_1}$ its
part of degree at least two has nonnegative coefficients, and
$\rho_a+\rho_c\le16q$, $\rho_a\rho_c\le(\rho_a+\rho_c)^2/4$, $\rho_{b_l}\le32q$;
so it is at most its value at $\rho_a=\rho_c=8q$, $\rho_{b_l}=32q$, which is
$(1+32q)^2F(8q,8q)-1-96q=P_c-96q$. All these bounds are pointwise in $x$, and
\cref{cor:incident-weights} turns them into the norm bound
$\norm{T_{\rm coll}}\le D_{\rm bal}$.

(b) \emph{Endpoint monotonicity.} Write every quantity as a function of
$r=p^{-1/2}$, with $q=r^2$, on $(0,P^{-1/2}]$. The functions $N_{\rm ns}/q$ and
$N_A/q$ are admissible in the sense of \cref{lem:majorant}
(\cref{cor:extraction-compressed,cor:extraction-leaf}), and so is
$D_{\rm bal}/q$: its numerator divided by $q$ is a polynomial with nonnegative
coefficients in $r$ and the count errors, and $\Lambda_c$ is a product of
factors $1-\varphi$ with admissible $\varphi<1$. The functions $m_\star/q$,
$\mathfrak m_+/q$ and $\chi^+_p/q$ are nondecreasing: written in $r$, the
numerator of $m_\star/q$ is nondecreasing on $(0,1]$ and its denominator $Z_4$
is positive and nonincreasing, and $\chi^+_p/q=(D_{\rm full}/r)^2/L_{\rm low}^2$
with $L_{\rm low}$ positive and nonincreasing (\cref{lem:record:root-marginal}(iii))
and $D_{\rm full}/r$ nondecreasing, since each term of $D_{\rm full}$ carries one
of the factors $\Delta_2$, $q$, $\mathfrak a_p$, and $\Delta_2/r$, $q/r$ and
$\mathfrak a_p/r$ are nondecreasing. The quantities $g_{\mathbf z}$,
$1/g_{\mathbf v}=1+q/(1-2q)$, $1/h_p=(1+3q/(1-4q))/(1-5q)$, $\rho/q$, $\eta/r$,
$s_p$ and $k_4$ are built from these by sums, products, square roots and
quotients by positive nonincreasing denominators. Hence $\eps_Z$, $H_{20}$,
$\tau_g/r$, $w_g/q$, $x_g/q=2r\,w_g/q+r(\tau_g/r)(N_A/q)$ and
$c^+_{20}=(\chi^+_p/q)/(1-H_{20})$ are nondecreasing in $r$ while $\eps_Z$ and
$H_{20}$ are below one; $\Delta_g/q=30-(l_g/r)^2$, with $l_g/r=2+r\,N_A/q$, is
nonincreasing; and where $\Delta_g>0$ the functions
$(c_{\rm gap}/r)^2=(\tau_g/r)^2+(x_g/q)^2/(\Delta_g/q)$ and
$C_{20}/r=(s_p(c_{\rm gap}/r)/h_p+D_{\rm bal}/r)/(1-\eps_Z)$ are nondecreasing.
If the gates hold at $P$, they therefore hold at every prime $p\ge P$, with
$g_{20}(p)\le g_{20}(P)$ and $c^+_{20}(p)\le c^+_{20}(P)$.

\emph{The caps.} At such a prime, \cref{prop:l20-pair-gap} gives
$\norm{\mathsf R_c}\le c_{\rm gap}<1$, \cref{prop:l20-rectangular}(d) with
$c_{\rm R}=c_{\rm gap}$ and $D_{\rm col}=D_{\rm bal}$ from (a) gives
$T\succeq-C_{20}I$ and $\norm{T|_{\mathcal I^\perp}}\le C_{20}$, and
\cref{prop:l20-completion} with $C_\perp=C_{20}$ gives
$-g_{20}(p)rI\preceq\mathsf A_p\preceq(1+c^+_{20}(p)/p)I$. The block bound follows
from \cref{lem:spec:tensorcomp}(b) with $g=g_{20}(P)$ and $c^+=c^+_{20}(P)$,
where $g_{20}(P)p^{-1/2}\le C_{20}(P)<1$, and from
$\sum_{p\in\Blk}1/p\le L/(X\log X)\le L_{\rm loc}/(X\log X)$
(\cref{lem:record:local-domain}(ii)).
\end{proof}

At the endpoint $p=Q_3$ (just below the middle primes) the main terms of $C_{20}/r$
are $s_pc_{\rm gap}/(h_pr)=5.47787\ldots$, against their limit
$\sqrt{30}=5.47723\ldots$ as $p\to\infty$, and
$D_{\rm bal}/r=0.02176\ldots$ ($D_{\rm bal}/q=97.18\ldots$,
close to its leading coefficient $96$); the normalizer
$h_p^{-1}$ adds about $4.01\cdot10^{-7}$ as a relative correction. So the
collision bound $D_{\rm bal}$ is the largest correction, and
$g_{20}-\sqrt{30}$ is asymptotic to $96r$ as $p\to\infty$. The
other inputs at $Q_3$ are $N_{\rm ns}=7.336\cdot10^{-5}$ up to
the displayed digits,
$N_A/q=281.37$, $k_4=5.4772264$, $c_{\rm gap}=1.226351\cdot10^{-3}$,
$\chi^+_p=1.0341\cdot10^{-6}$ and $H_{20}=7.8240\cdot10^{-5}$.

\begin{certificate}[The L20 grid]\label{cert:l20-grid}
At each of the $45$ endpoints below, the quantities
$N_{\rm ns}$, $N_A$, $D_{\rm bal}$, $m_\star$, $\eps_Z$, $k_4$,
$s_P$, $h_P$, $c_{\rm gap}$, $C_{20}$, $H_{20}$ and $\chi^+_P$ of
the L20 estimates satisfy $\eps_Z<1/2$, $\Delta_g>0$,
$c_{\rm gap}<1$, $C_{20}<1$ and $H_{20}<1$ in outward rational
interval arithmetic, with the recorded values of $g_{20}$ and
$c^+_{20}$ stated below. The scripts and records are in the
directory \path{research/advance_2026_10_24/local/}. The record
\path{l20_refined_Q_grid.json}, built by \path{l20_collision_balanced.py} on the
rank-four, rectangular and gap evaluators \path{l20_rank4.py},
\path{l20_rectangular.py} and \path{l20_rectangular_gap.py} (and on the star
matrix and the frame cap of \path{l20_full_star.py}, which asserts the
characteristic polynomial of \cref{lem:l20-star-matrix}(a) in exact symbolic
arithmetic), contains $45$ endpoints from $2^{20}$ to $2^{50}$, among them the
$33$ piece starts. The nonstar ledger of $N_A$ (\cref{prop:l20-rectangular}(b))
is the histogram \path{middle_factor_probe/rectangular_A.hist.json.gz} of the
same directory. At each endpoint $P$ it evaluates in outward rational
interval arithmetic $N_{\rm ns}$, $N_A$, $D_{\rm bal}$, $m_\star$, $\eps_Z$,
$k_4$, $s_P$, $h_P$, $c_{\rm gap}$, $C_{20}$, $H_{20}$ and $\chi^+_P$, checks
$\eps_Z<1/2$, $\Delta_g>0$, $c_{\rm gap}<1$, $C_{20}<1$ and $H_{20}<1$, and
records $g_{20}(P)$, rounded up to $10^{-12}$, and $c^+_{20}(P)$. At $Q_3$ it
gives $g_{20}=5.500057067029$ and $c^+_{20}<20.633953751661$, hence a block
cap below $5.500255036718$ at $L_{\rm loc}$; at $2^{20}$, $g_{20}=5.597878403332$;
the limit is $\sqrt{30}$. These are the inherited uniform-root prices, preserved as a baseline in
\cref{tab:data:l20-grid}. The selected native reconstruction recomputes
the protected collision price of \cref{imp:local:collision-gap} at all
33 piece starts, with the same positive coefficient and marginal data.
It then replaces L20 and L21 on the seven middle pieces from $2^{30}$ by the
prices of \cref{cert:positive-root-caps}. The reference uses
the selected value at $Q_3$, $g_{20}=5.481634403216$, proved
through \cref{imp:local:collision-gap}; the larger value
$5.500057067029$ recorded above is the weaker bound of
\cref{prop:l20-collisions}, retained only as a comparison. At
$Q_3$ the selected uniform root also has
$c^+_{20}=20.633953751661$ and
block cap below $5.481831709798$ at budget $585$.
The exact selected spectra and weights are in
\path{research/push_2026_10_25/global/refined_L570_source.json};
\path{research/improvements_2026_10_25/verify.py} independently
checks the collision, endpoint and native-completion interfaces,
and \path{research/push_2026_10_25/verify.py} authenticates their
unchanged identity after scalar reanchoring.
\end{certificate}

\paragraph{Consequences and what to read next.}
\Cref{thm:bmor} (the cell charges above $X_{\rm Ch}$); \cref{lem:paley-triple}
(class-three counts); \cref{lem:record:fold-identification},
\cref{def:graphs:labeling} and \cref{def:spec:data} (the twelve-label graphs
and their cuts); \cref{thm:spectral-criterion} and
\cref{cor:spectral-criterion-blocks}, with the inputs of
\cref{cor:extraction-compressed,lem:oct24-collision} (the route-A caps);
\cref{lem:majorant} (evaluation at the piece start and at the endpoints);
\cref{lem:spec:expansion,lem:spec:stars}, \cref{lem:general-star},
\cref{cor:star-matrix}, \cref{lem:spec:affpair}(d), \cref{lem:l20-star-mean} and
\cref{lem:record:weighted-moment}, \cref{lem:swapped-rectangle},
\cref{prop:matching:positive}(a$''$), \cref{thm:weil}, \eqref{eq:two-root},
\cref{cor:soft-kernels}(a), \cref{cor:extraction-leaf}, \cref{cor:incident-weights} and
\cref{lem:record:root-marginal}(ii),(iii) (the cap of L20);
\cref{lem:spec:cutkernel} and \cref{lem:spec:tensorcomp}(b) (block caps);
\cref{lem:record:repricing} and \cref{def:closing-quantities} (prices and the
selection factor at $a$). Certificates:
\cref{cert:record:middle-price,cert:layer12-grid,cert:l20-grid}.

\paragraph{Hypotheses verified in this section (\cref{tab:hypotheses}).} the inputs of the rows ``(CP1)
middle'' and ``$P$, $V$, $D$, $U_0$, tags'' that do not depend on the reference:
the middle partition and its charges
(\cref{def:record:exact-blocks,lem:record:block-charge,lem:record:lower-start-sums},
\cref{cert:record:middle-price}) and the caps of the twelve-label cuts L21, E21
and L10 and of the cut L20 of the uniform root
(\cref{lem:layer12-caps,prop:l20-collisions};
\cref{cert:layer12-grid,cert:l20-grid}); and the piece caps of the restorations
(\cref{lem:piece-caps}), an input of the row ``(CP5) variance classes''. The
first two rows are completed in \cref{sec:middle-frame}.
